\part{Datos y Patrones Arquitectónicos}

\chapter{Gestión de estado y objetos integrados en aplicaciones web}
\label{cap:semana10}

\epigraph{«HTTP es el protocolo de la amnesia colectiva: cada petición
nace y muere sin memoria. Programar la web es, en buena medida, el arte
de simular una memoria que el protocolo nos negó.»}{--- \textit{Adaptado
de Roy Fielding}}

\section*{Resultado de aprendizaje de la semana}
Al concluir la semana 10, el estudiante diseña e implementa soluciones
web que integran mecanismos robustos de gestión de estado --- cookies,
sesiones de servidor, almacenamiento del cliente y tokens --- y domina
los objetos integrados \texttt{Request}, \texttt{Response},
\texttt{Session}, \texttt{Application} y \texttt{Page} en PHP, Java y
ASP.NET Core, aplicando los flags de seguridad estándar que la
industria exige en 2026.

\section{Introducción: el dilema fundamental de HTTP}

HTTP es un protocolo \emph{stateless}: cada petición se
trata independientemente, sin recordar las anteriores. Esto es una
fortaleza para escalar (cualquier servidor puede atender cualquier
petición), pero un problema para construir aplicaciones que necesitan
recordar al usuario entre clics. La resolución de esta tensión es el
tema central de la semana.


\begin{cajahistoria}{HTTP --- el protocolo amnésico que cambió el mundo}
Cuando Tim Berners-Lee diseñó HTTP en 1989, lo hizo deliberadamente
como un protocolo \emph{stateless} (sin estado): cada petición es
independiente, autocontenida, y el servidor no recuerda nada de la
petición anterior. Esta decisión, tomada en una época de servidores
con poca memoria y conexiones lentas, parecía limitada pero resultó
genial: permitió que la web escalara a miles de millones de usuarios
sin colapsar bajo el peso de la memoria de sesiones.

Pero la web pronto necesitó \emph{algo} que pareciera memoria: ¿cómo
saber qué usuario está autenticado?, ¿cómo recordar los productos en
el carrito?, ¿cómo personalizar la experiencia? La respuesta llegó en
1994 cuando Lou Montulli, ingeniero de Netscape, inventó las
\textbf{cookies HTTP} originalmente como solución para un cliente de
e-commerce. Las cookies fueron una idea simple: el servidor envía un
identificador en la respuesta; el navegador se lo guarda; en cada
petición posterior, el navegador devuelve ese identificador. Con eso,
el servidor puede asociar peticiones del mismo usuario.

Sobre las cookies se construyeron las \textbf{sesiones de servidor}:
el servidor guarda los datos del usuario en su propia memoria
asociados a un \emph{session ID} aleatorio, y solo el ID viaja en la
cookie. Treinta años después, este modelo sigue siendo el sustrato de
la mayoría de aplicaciones web tradicionales del mundo. Pero la
arquitectura moderna --- microservicios, escalado horizontal, SPAs
--- llevó a la emergencia de los \textbf{tokens stateless} (JWT) que
trasladan el estado de vuelta al cliente, permitiendo que el servidor
permanezca verdaderamente sin memoria.

La literatura técnica reciente confirma que la gestión de estado sigue
siendo uno de los aspectos más críticos de la arquitectura web:
\cite{fielding2000rest} en su tesis doctoral fundacional estableció
que \emph{statelessness} es una restricción arquitectónica deseable
para escalabilidad; \cite{kleppmann2017designing} discute en
profundidad cómo el estado distribuido es la fuente de la mayoría de
los bugs sutiles en sistemas modernos; y
\cite{jones2015jwt} (RFC 7519) estandarizó JWT como el mecanismo
dominante para sesiones distribuidas, hoy ampliamente adoptado en
APIs públicas \cite{jones2015jwt}.
\end{cajahistoria}

\subsection{Las cinco estrategias de gestión de estado}

Hay cinco estrategias diferentes para preservar estado entre
peticiones HTTP. Cada una tiene ventajas, desventajas y casos donde
es la mejor opción. La tabla siguiente las contrasta.


\begin{table}[htbp]
\centering
\caption{Estrategias de gestión de estado en aplicaciones web modernas.}
\footnotesize
\begin{tabularx}{\textwidth}{lXXl}
\toprule
\textbf{Mecanismo} & \textbf{Ubicación del estado} & \textbf{Caso de uso ideal} & \textbf{Tamaño}\\
\midrule
Cookie clásica & Cliente (cookie HTTP) & Preferencias simples (idioma, tema) & 4 KB\\
Sesión servidor & Servidor (memoria/Redis) & App tradicional multi-página & MB \\
\texttt{localStorage} & Cliente (navegador) & Datos persistentes locales & 5--10 MB\\
\texttt{sessionStorage} & Cliente (pestaña) & Datos volátiles de pestaña & 5--10 MB\\
JWT & Cliente (cookie/header) & APIs REST stateless, SPAs & 1--4 KB\\
IndexedDB & Cliente (navegador) & Apps offline, gran volumen & GB\\
\bottomrule
\end{tabularx}
\end{table}

\subsection{Los flags de seguridad de cookies en 2026}

Antes de profundizar en cualquier mecanismo, hay que dominar los
\emph{flags} de seguridad que toda cookie debe llevar en producción:

\begin{cajadefinicion}{Los cuatro flags imprescindibles}
\begin{description}[leftmargin=1.5cm]
\item[\texttt{Secure}.] La cookie solo viaja por HTTPS, nunca por HTTP
plano. \emph{Sin este flag, la cookie puede ser interceptada en redes
WiFi públicas.}
\item[\texttt{HttpOnly}.] JavaScript del navegador no puede leer la
cookie a través de \texttt{document.cookie}. \emph{Defensa principal
contra XSS para robar la sesión.}
\item[\texttt{SameSite}.] Controla cuándo el navegador envía la cookie
en peticiones de otros sitios. Tres valores: \texttt{Strict} (nunca
desde otros sitios), \texttt{Lax} (solo en navegación de nivel
superior), \texttt{None} (siempre, requiere \texttt{Secure}).
\emph{Defensa principal contra CSRF.}
\item[\texttt{Max-Age} o \texttt{Expires}.] Tiempo de vida de la
cookie. Sin esto, es \emph{session cookie} (se borra al cerrar el
navegador).
\end{description}
\end{cajadefinicion}

\begin{cajaadvertencia}{Combinación correcta vs incorrecta}
\textbf{Correcto para sesiones de autenticación (Listado~\ref{lst:10.1}):}
\begin{lstlisting}[style=shell,numbers=none,caption={Ejemplo de Shell},label=lst:10.1]
Set-Cookie: SESSID=abc123; Path=/; Secure; HttpOnly; SameSite=Strict; Max-Age=1800
\end{lstlisting}
\textbf{Incorrecto (frecuentemente visto en código de aprendizaje) (Listado~\ref{lst:10.2}):}
\begin{lstlisting}[style=shell,numbers=none,caption={Ejemplo de Shell},label=lst:10.2]
Set-Cookie: SESSID=abc123
\end{lstlisting}
La segunda forma expone la cookie a HTTP plano (puede leerse en redes
inseguras), permite que JavaScript la lea (vulnerabilidad XSS), y se
envía en peticiones cross-site (vulnerabilidad CSRF).
\end{cajaadvertencia}

\section{El objeto Request: lectura segura de la petición}

El objeto \texttt{Request} encapsula toda la información que el
cliente envía al servidor. Su correcta lectura es la primera línea de
defensa de la aplicación: nunca se debe confiar en datos del cliente
sin filtrarlos.

\subsection{Componentes del objeto Request}

El objeto \emph{Request} encapsula todo lo que el cliente
envía: URL, método, cabeceras, cuerpo, cookies. Conocer sus
componentes es prerrequisito para programar cualquier endpoint
profesional.

\textbf{Ejemplo corto comentado (Listado~\ref{lst:10.3}):}
\begin{lstlisting}[language=PHP,numbers=none,caption={Ejemplo de PHP},label=lst:10.3]
<?php
// Componentes principales del Request HTTP en PHP
$_SERVER['REQUEST_METHOD']; // 'GET', 'POST', 'PUT', 'DELETE'
$_SERVER['REQUEST_URI']; // '/api/v1/usuarios?id=42'
$_SERVER['HTTP_HOST']; // 'tienda.uteq.edu.ec'
$_SERVER['HTTP_USER_AGENT']; // navegador del cliente
$_GET['id']; // parametros de la URL
$_POST['email']; // datos del cuerpo (form-urlencoded)
$_COOKIE['SESSID']; // cookies enviadas por el cliente
file_get_contents('php://input'); // cuerpo crudo (JSON, XML)
?>
\end{lstlisting}


\begin{itemize}
\item \textbf{Método HTTP}: \texttt{GET}, \texttt{POST}, \texttt{PUT},
\texttt{DELETE}, \texttt{PATCH}.
\item \textbf{URL completa}: ruta, query string, fragmento.
\item \textbf{Encabezados}: \texttt{Content-Type},
\texttt{Authorization}, \texttt{User-Agent}, \texttt{Accept-Language}.
\item \textbf{Cookies}: enviadas por el navegador.
\item \textbf{Cuerpo} (en POST/PUT/PATCH): JSON, form-urlencoded,
multipart, raw.
\item \textbf{Metadatos del cliente}: IP origen, protocolo (HTTP/HTTPS).
\end{itemize}

\begin{cajaejemplo}{Ejemplo 10.1 --- Lectura segura del Request en PHP 8.4}

\textbf{Enunciado.} Mostrar la lectura segura del objeto Request en PHP 8.4 obteniendo método, URL, cabeceras, query string y cuerpo, con validación básica de cada campo.

\begin{lstlisting}[style=php,caption={request-helpers.php},label=lst:10.4]
<?php
declare(strict_types=1);

/**
 * Obtiene la IP real del cliente, considerando proxy reverso.
 * Verifica las cabeceras X-Forwarded-For y X-Real-IP que un proxy
 * Nginx o un balanceador de carga (HAProxy, AWS ELB) inyecta.
 */
function obtenerIpCliente(): string {
    $cabeceras = [
        'HTTP_X_FORWARDED_FOR', // proxy estandar
        'HTTP_X_REAL_IP', // Nginx
        'HTTP_CF_CONNECTING_IP', // Cloudflare
        'REMOTE_ADDR' // conexion directa
    ];
    foreach ($cabeceras as $cabecera) { // itera sobre colecion
        if (!empty($_SERVER[$cabecera])) { // condicional
            // X-Forwarded-For puede tener varias IPs separadas por coma
            $ip = explode(',', $_SERVER[$cabecera])[0];
            $ip = trim($ip);
            if (filter_var($ip, FILTER_VALIDATE_IP, // valida/sanea entrada con filtro
                FILTER_FLAG_NO_PRIV_RANGE | FILTER_FLAG_NO_RES_RANGE)) {
                return $ip; // retorna valor
            }
        }
    }
    return '0.0.0.0'; // retorna valor
}

/**
 * Lectura segura de un parametro GET.
 * NUNCA acceder directamente a $_GET['param'] sin filtrar.
 */
function leerEntero(string $nombre, int $minimo = 1,
                    int $maximo = PHP_INT_MAX, int $defecto = 1): int {
    $val = filter_input(INPUT_GET, $nombre, FILTER_VALIDATE_INT,
        ['options' => [
            'default'   => $defecto,
            'min_range' => $minimo,
            'max_range' => $maximo
        ]]);
    return $val ?? $defecto; // retorna valor
}

function leerCadena(string $nombre, int $maxLen = 100,
                    string $defecto = ''): string {
    $raw = $_REQUEST[$nombre] ?? $defecto;
    $val = trim((string)$raw);
    // condicional
    // condicional
    if (strlen($val) > $maxLen) $val = substr($val, 0, $maxLen);
    return $val; // retorna valor
}

function exigirMetodo(string $metodo): void {
    // condicional
    // condicional
    if ($_SERVER['REQUEST_METHOD'] !== strtoupper($metodo)) {
        http_response_code(405);
        header('Allow: ' . strtoupper($metodo)); // envia cabecera HTTP al cliente
        exit("Metodo no permitido"); // termina ejecucion
    }
}

// Uso conjunto
exigirMetodo('GET');
$pagina  = leerEntero('page', 1, 1000, 1);
$busca   = leerCadena('q', 80);
$ip      = obtenerIpCliente();

// Registrar acceso para auditoria
error_log("[" . date('c') . "] $ip - GET ?page=$pagina&q=$busca");
\end{lstlisting}

\textbf{Por qué este código es importante:}
\begin{itemize}
\item Centraliza la lectura segura en funciones helper reutilizables.
\item Detecta proxy reverso correctamente (un error común en
producción).
\item Filtra IPs privadas y reservadas con flags de
\texttt{filter\_var}.
\item Aplica límites estrictos a los enteros (rangos).
\item Trunca cadenas demasiado largas para prevenir ataques DoS por
memoria.
\item Verifica el método HTTP esperado.
\end{itemize}
\end{cajaejemplo}

\subsection{Lectura del Request en Java Servlets / Spring Boot}

En Java/Spring Boot, el equivalente al \texttt{\$\_SERVER}
de PHP es la interfaz \texttt{HttpServletRequest} o las anotaciones
de Spring (\texttt{@RequestParam}, \texttt{@RequestHeader},
\texttt{@RequestBody}). El ejemplo siguiente las recorre.


\begin{cajaejemplo}{Ejemplo 10.2 --- Acceso al Request en Spring Boot 4}

\textbf{Enunciado.} Mostrar el acceso al objeto Request en Spring Boot 4 usando las anotaciones idiomáticas: \texttt{@RequestParam}, \texttt{@RequestHeader}, \texttt{@RequestBody}.

\begin{lstlisting}[style=java,caption={RequestInspectorController.java},label=lst:10.5]
package ec.edu.uteq.appweb.controller; // paquete del archivo

import jakarta.servlet.http.HttpServletRequest; // importacion de clase externa
import org.springframework.web.bind.annotation.*; // importacion de clase externa
import java.util.*; // importacion de clase externa

@RestController // expone endpoints REST con JSON
@RequestMapping("/api/inspeccion") // ruta base del controlador
public class RequestInspectorController { // clase publica

    @GetMapping // endpoint HTTP GET
    public Map<String, Object> inspeccionar( // metodo publico
            HttpServletRequest req,
            @RequestParam(defaultValue = "") String q, // extrae parametro de query string
            @RequestParam(defaultValue = "1") int page, // extrae parametro de query string
            @RequestHeader(value = "User-Agent", // extrae cabecera HTTP
                           defaultValue = "?") String userAgent) {

        // Headers selectivos
        var headers = new LinkedHashMap<String, String>();
        var nombres = req.getHeaderNames();
        while (nombres.hasMoreElements()) { // bucle while
            var n = nombres.nextElement();
            headers.put(n, req.getHeader(n));
        }

        return Map.of( // retorna
            "metodo",      req.getMethod(),
            "url",         req.getRequestURL().toString(),
            "remoteAddr",  obtenerIp(req),
            "userAgent",   userAgent,
            "queryParams", Map.of("q", q, "page", page),
            "headersSize", headers.size()
        );
    }

    private String obtenerIp(HttpServletRequest r) { // campo privado
        var fwd = r.getHeader("X-Forwarded-For");
        if (fwd != null && !fwd.isBlank()) { // condicional
            return fwd.split(",")[0].trim(); // retorna
        }
        return r.getRemoteAddr(); // retorna
    }
}
\end{lstlisting}
\end{cajaejemplo}

\subsection{Lectura del Request en ASP.NET Core 10}

ASP.NET Core usa \texttt{HttpContext.Request} para acceder a
los datos de la petición, con una API moderna basada en propiedades
fuertemente tipadas. El ejemplo siguiente la ilustra.


\begin{cajaejemplo}{Ejemplo 10.3 --- Endpoint que inspecciona el Request en ASP.NET Core}

\textbf{Enunciado.} Implementar un endpoint en ASP.NET Core que inspeccione el Request: método, ruta, cabeceras, query string. Se usará el objeto \texttt{HttpContext.Request}.

\begin{lstlisting}[style=cs,caption={Controllers/InspeccionController.cs},label=lst:10.6]
using Microsoft.AspNetCore.Mvc; // importa namespace

namespace AppwebApi.Controllers; // namespace del archivo

[ApiController] // controlador de Web API (validacion automatica)
[Route("api/[controller]")] // ruta base del controlador
public class InspeccionController : ControllerBase // clase publica
{
    [HttpGet] // endpoint HTTP GET
    public IActionResult Inspeccionar( // metodo o miembro publico
        [FromQuery] string q = "", // mapea query string
        [FromQuery] int page = 1) // mapea query string
    {
        // IP real considerando X-Forwarded-For
        var ip = HttpContext.Request.Headers["X-Forwarded-For"]
                    .FirstOrDefault()?.Split(',')[0].Trim()
                 ?? HttpContext.Connection.RemoteIpAddress?.ToString()
                 ?? "0.0.0.0";

        return Ok(new { // HTTP 200 con cuerpo
            metodo      = Request.Method,
            url         = $"{Request.Scheme}://{Request.Host}{Request.Path}",
            remoteAddr  = ip,
            userAgent   = Request.Headers.UserAgent.ToString(),
            queryParams = new { q, page },
            cookieCount = Request.Cookies.Count
        });
    }
}
\end{lstlisting}
\end{cajaejemplo}

\section{El objeto Response: control total de la respuesta HTTP}

El objeto \texttt{Response} permite controlar tres aspectos:
código de estado, cabeceras y cuerpo.

\subsection{Códigos de estado HTTP --- la guía exhaustiva}

Los códigos de estado HTTP son la \emph{lengua común} entre
cliente y servidor para indicar el resultado de una petición.
Confundirlos (devolver 200 cuando hubo un error, por ejemplo) es uno
de los defectos más comunes en APIs amateur.


\begin{table}[htbp]
\centering
\caption{Códigos de estado HTTP de uso obligado en aplicaciones web profesionales.}
\footnotesize
\begin{tabularx}{\textwidth}{llX}
\toprule
\textbf{Código} & \textbf{Nombre} & \textbf{Cuándo usarlo (con ejemplos)}\\
\midrule
\multicolumn{3}{l}{\textbf{2xx --- Éxito}}\\
200 & OK                  & Operación exitosa con cuerpo en respuesta. GET, PUT exitosos.\\
201 & Created             & Recurso creado. POST exitoso. Incluir \texttt{Location}.\\
204 & No Content          & Operación exitosa sin cuerpo. DELETE típicamente.\\
\midrule
\multicolumn{3}{l}{\textbf{3xx --- Redirección}}\\
301 & Moved Permanently   & Cambio definitivo de URL. SEO sigue al nuevo URL.\\
302 & Found               & Redirección temporal. POST $\to$ GET (patrón PRG).\\
304 & Not Modified        & Caché válida (responde a \texttt{If-None-Match}).\\
\midrule
\multicolumn{3}{l}{\textbf{4xx --- Error del cliente}}\\
400 & Bad Request         & JSON mal formado, parámetro faltante crítico.\\
401 & Unauthorized        & No autenticado. Falta JWT o cookie de sesión.\\
403 & Forbidden           & Autenticado pero sin permiso. Rol insuficiente.\\
404 & Not Found           & Recurso inexistente. ID inválido.\\
405 & Method Not Allowed  & Endpoint existe pero no acepta ese verbo. Incluir \texttt{Allow}.\\
409 & Conflict            & Email duplicado, optimistic locking failure.\\
422 & Unprocessable Entity & Validación de negocio falla (campos inválidos).\\
429 & Too Many Requests   & Rate limit excedido. Incluir \texttt{Retry-After}.\\
\midrule
\multicolumn{3}{l}{\textbf{5xx --- Error del servidor}}\\
500 & Internal Server Error & Excepción no controlada. NUNCA mostrar stack trace al cliente.\\
502 & Bad Gateway         & Proxy recibió respuesta inválida del upstream.\\
503 & Service Unavailable & Mantenimiento programado. Incluir \texttt{Retry-After}.\\
504 & Gateway Timeout     & Upstream tardó demasiado.\\
\bottomrule
\end{tabularx}
\end{table}

\subsection{Errores estandarizados con RFC 7807 ProblemDetails}

El RFC 7807 \cite{rfc7807} estandariza un formato JSON para errores
HTTP. Es el estándar de la industria en 2026 y los frameworks modernos
(ASP.NET Core, Spring Boot) lo soportan nativamente (Listado~\ref{lst:10.7}):

\begin{lstlisting}[style=js,numbers=none,caption={Respuesta de error según RFC 7807},label=lst:10.7]
HTTP/1.1 422 Unprocessable Entity
Content-Type: application/problem+json

{
  "type":   "https://uteq.edu.ec/errors/validacion",
  "title":  "Datos invalidos",
  "status": 422,
  "detail": "El campo email no cumple el formato esperado",
  "instance": "/api/usuarios",
  "errores": {
    "email":  ["debe tener formato valido"],
    "edad":   ["debe ser mayor a 0"]
  }
}
\end{lstlisting}

\subsection{Cabeceras de seguridad obligatorias en producción}

En producción, además de las cabeceras de seguridad básicas
vistas en la semana 9, conviene configurar otras que protegen contra
ataques específicos. La tabla y los ejemplos siguientes las catalogan.


\begin{cajabuenas}{Las siete cabeceras de seguridad imprescindibles}
\begin{enumerate}
\item \textbf{Strict-Transport-Security}:
\texttt{max-age=31536000; includeSubDomains; preload}\\
Fuerza HTTPS por un año en todo el dominio.

\item \textbf{Content-Security-Policy}:
\texttt{default-src 'self'; script-src 'self' 'nonce-XYZ'}\\
Mitiga XSS limitando fuentes de scripts y estilos.

\item \textbf{X-Frame-Options}: \texttt{DENY}\\
No permite que el sitio se cargue en \texttt{<iframe>}: anti-clickjacking.

\item \textbf{X-Content-Type-Options}: \texttt{nosniff}\\
Impide que el navegador deduzca el tipo MIME (anti-MIME-sniffing).

\item \textbf{Referrer-Policy}:
\texttt{strict-origin-when-cross-origin}\\
Equilibrio entre privacidad y funcionalidad.

\item \textbf{Permissions-Policy}:
\texttt{geolocation=(), microphone=(), camera=()}\\
Restringe APIs sensibles del navegador.

\item \textbf{Cross-Origin-Opener-Policy}: \texttt{same-origin}\\
Aísla el contexto de navegación de otros orígenes (mitigación
Spectre).
\end{enumerate}
\end{cajabuenas}

\begin{cajaejemplo}{Ejemplo 10.4 --- Cabeceras de seguridad en las tres plataformas}

\textbf{Enunciado.} Aplicar las cabeceras de seguridad obligatorias en producción (\texttt{Content-Security-Policy}, \texttt{X-Frame-Options}, \texttt{Strict-Transport-Security}, etc.) en las tres plataformas del curso.

\textbf{En PHP (Listado~\ref{lst:10.8}):}
\begin{lstlisting}[style=php,numbers=none,caption={Ejemplo de PHP},label=lst:10.8]
<?php
// envia cabecera HTTP al cliente
// envia cabecera HTTP al cliente
header('Strict-Transport-Security: max-age=31536000; includeSubDomains');
header("Content-Security-Policy: default-src 'self'"); // envia cabecera HTTP al cliente
header('X-Frame-Options: DENY'); // envia cabecera HTTP al cliente
header('X-Content-Type-Options: nosniff'); // envia cabecera HTTP al cliente
header('Referrer-Policy: strict-origin-when-cross-origin'); // envia cabecera HTTP al cliente
// envia cabecera HTTP al cliente
// envia cabecera HTTP al cliente
header('Permissions-Policy: geolocation=(), microphone=(), camera=()');
\end{lstlisting}

\textbf{En Spring Boot 4 (Listado~\ref{lst:10.9}):}
\begin{lstlisting}[style=java,numbers=none,caption={Ejemplo de Java},label=lst:10.9]
@Configuration // clase de configuracion Spring
public class SeguridadConfig { // clase publica
    @Bean
    // metodo publico
    // metodo publico
    public SecurityFilterChain seguridad(HttpSecurity http) throws Exception {
        http.headers(h -> h
            .httpStrictTransportSecurity(hsts -> hsts
                .maxAgeInSeconds(31536000).includeSubDomains(true))
            .frameOptions(f -> f.deny())
            .contentSecurityPolicy(csp -> csp
                .policyDirectives("default-src 'self'"))
            .referrerPolicy(rp -> rp.policy(
                ReferrerPolicy.STRICT_ORIGIN_WHEN_CROSS_ORIGIN))
        );
        return http.build(); // retorna
    }
}
\end{lstlisting}

\textbf{En ASP.NET Core 10 (Listado~\ref{lst:10.10}):}
\begin{lstlisting}[style=cs,numbers=none,caption={Ejemplo de C\#},label=lst:10.10]
app.Use(async (ctx, next) => {
    ctx.Response.Headers["Strict-Transport-Security"]
        = "max-age=31536000; includeSubDomains";
    ctx.Response.Headers["X-Frame-Options"] = "DENY";
    ctx.Response.Headers["X-Content-Type-Options"] = "nosniff";
    ctx.Response.Headers["Referrer-Policy"]
        = "strict-origin-when-cross-origin";
    ctx.Response.Headers["Content-Security-Policy"]
        = "default-src 'self'";
    await next();
});
\end{lstlisting}
\end{cajaejemplo}

\section{El objeto Session: estado entre peticiones}

La \emph{sesión} es el mecanismo más antiguo y aún
ampliamente usado para mantener estado entre peticiones: el servidor
guarda los datos del usuario en memoria (o en Redis) y envía al
cliente solo un identificador.


\subsection{Mecánica detallada de las sesiones de servidor}

Comprender la mecánica detallada de las sesiones de servidor
es clave porque la mayoría de fallos de seguridad ocurren por
desconocimiento de cómo se genera, transporta y persiste el
\emph{Session ID}.


\begin{figure}[htbp]
\centering
\begin{tikzpicture}[
  node distance=1.2cm and 1.5cm,
  every node/.style={font=\scriptsize,align=center},
  caja/.style={rectangle,draw=uteqverde,thick,minimum width=2cm,
               minimum height=0.8cm,fill=uteqverde!10,rounded corners=2pt},
  flecha/.style={->,>=stealth,thick}]

  \node[caja] (cli)  {Navegador};
  \node[caja,right=4cm of cli] (srv) {Servidor};
  \node[cylinder,draw=uteqverde,fill=uteqverde!15,thick,
        shape border rotate=90,aspect=0.3,minimum width=1.5cm,
        minimum height=1.7cm,below=0.5cm of srv] (sto)
        {Session\\Store};

  \draw[flecha] ([yshift=0.4cm]cli.east) --
    node[above]{1. GET /login (sin cookie)} ([yshift=0.4cm]srv.west);
  \draw[flecha] ([yshift=-0.4cm]srv.west) --
    node[below,align=center]{\scriptsize 2. Set-Cookie: SESSID=abc...\\\scriptsize Secure; HttpOnly; SameSite}
    ([yshift=-0.4cm]cli.east);
  \draw[flecha] (srv) -- node[right]{3. Crear} (sto);

\end{tikzpicture}
\caption{Establecimiento de una sesión: el servidor genera un ID
aleatorio, lo asocia con datos en su \emph{Session Store} (memoria,
Redis o BD) y lo envía al cliente como cookie con flags de seguridad.}
\end{figure}

\subsection{Implementación comparada en las cuatro plataformas}

Cada plataforma implementa sesiones con su propia API, pero los
conceptos son los mismos. La tabla siguiente las contrasta en el orden
del libro ---Spring Boot primero---: Spring Boot, Django, PHP y
ASP.NET Core.


\begin{cajacompara}{Comparativa de APIs de sesión}
\centering
\scriptsize
\renewcommand{\arraystretch}{1.25}
\begin{tabularx}{\textwidth}{l
  >{\raggedright\arraybackslash}X
  >{\raggedright\arraybackslash}X
  >{\raggedright\arraybackslash}X
  >{\raggedright\arraybackslash}X}
\toprule
\textbf{Operación} & \textbf{Spring Boot 4} & \textbf{Django 5.2} & \textbf{PHP 8.4} & \textbf{ASP.NET Core 10}\\
\midrule
Iniciar      & automático (Servlet)       & \emph{middleware} automático & \texttt{session\_start()}        & \texttt{UseSession()}\\
Guardar      & \texttt{setAttribute(k,v)} & \texttt{session[k]=v}        & \texttt{\$\_SESSION[k]=\$v}      & \texttt{Session.SetString}\\
Leer         & \texttt{getAttribute(k)}   & \texttt{session[k]}          & \texttt{\$\_SESSION[k]}          & \texttt{Session.GetString}\\
Borrar       & \texttt{removeAttribute(k)}& \texttt{del session[k]}      & \texttt{unset(\$\_SESSION[k])}   & \texttt{Session.Remove}\\
Destruir     & \texttt{invalidate()}      & \texttt{session.flush()}     & \texttt{session\_destroy()}      & \texttt{Session.Clear()}\\
Regenerar ID & \texttt{changeSessionId()} & \texttt{cycle\_key()}        & \texttt{session\_\allowbreak regenerate\_id()} & regenerar cookie\\
Cookie       & \texttt{JSESSIONID}        & \texttt{sessionid}           & \texttt{PHPSESSID}               & \texttt{AspNetCore.Session}\\
Almacén      & JVM / Redis                & BD / cache / archivo         & archivos \texttt{/tmp}           & Memoria / Redis / SQL\\
Timeout      & 30 min                     & 2 semanas                    & 24 min                           & 20 min\\
\bottomrule
\end{tabularx}
\end{cajacompara}

\subsection{Sesiones en PHP 8.4 con configuración profesional}

PHP gestiona sesiones con la función nativa
\texttt{session\_start()}. El ejemplo siguiente la configura con todos
los flags de seguridad de producción y muestra los anti-patrones
comunes a evitar.


\begin{cajaejemplo}{Ejemplo 10.5 --- Sesión segura completa con login PHP}

\textbf{Enunciado.} Construir una sesión segura completa con login PHP: hash BCrypt, cookies con todos los flags de seguridad, regeneración de ID y logout robusto.

\begin{lstlisting}[style=php,caption={config/session.php (incluir antes de session\_start)},label=lst:10.11]
<?php
declare(strict_types=1);

/**
 * Configuracion de sesion para produccion.
 * DEBE llamarse ANTES de session_start().
 */
function configurarSesionSegura(): void {
    // Solo cookies HTTPS
    ini_set('session.cookie_secure',  '1');
    // JavaScript no puede leer la cookie
    ini_set('session.cookie_httponly', '1');
    // Anti-CSRF
    ini_set('session.cookie_samesite', 'Strict');
    // Forzar regeneracion de ID si llega uno desconocido
    ini_set('session.use_strict_mode', '1');
    // Timeout 30 min de inactividad
    ini_set('session.gc_maxlifetime',  '1800');
    // Cookie expira al cerrar navegador (=0) o tras 30 min (1800)
    session_set_cookie_params([ // flags de cookie
        'lifetime' => 1800,
        'path'     => '/',
        'secure'   => true,
        'httponly' => true,
        'samesite' => 'Strict',
    ]);
}
\end{lstlisting}

\begin{lstlisting}[style=php,caption={login.php con todas las prácticas},label=lst:10.12]
<?php
declare(strict_types=1);
require_once __DIR__ . '/config/session.php'; // incluye otro archivo PHP (fatal si falla)
require_once __DIR__ . '/config/db.php'; // incluye otro archivo PHP (fatal si falla)

configurarSesionSegura();
session_start(); // arranca o reanuda la sesion

// Si ya esta autenticado, redirigir
if (!empty($_SESSION['user_id'])) { // condicional
    header('Location: /dashboard.php'); // envia cabecera HTTP al cliente
    exit;
}

if ($_SERVER['REQUEST_METHOD'] === 'POST') { // condicional
    // Validar token CSRF (anti-CSRF en formularios)
    $tokenEnviado = $_POST['_csrf'] ?? '';
    $tokenSesion  = $_SESSION['_csrf'] ?? '';
    if (!hash_equals($tokenSesion, $tokenEnviado)) { // condicional
        http_response_code(419);
        exit('Token CSRF invalido'); // termina ejecucion
    }

    $email = filter_input(INPUT_POST, 'email',
                          FILTER_VALIDATE_EMAIL);
    $pass  = $_POST['password'] ?? '';

    if (!$email || strlen($pass) < 8) { // condicional
        http_response_code(422);
        echo 'Datos invalidos'; // imprime al output del servidor
        exit;
    }

    // Rate limiting en memoria (en produccion: Redis)
    $intentosKey = 'intentos:' . $_SERVER['REMOTE_ADDR'];
    if (($_SESSION[$intentosKey] ?? 0) >= 5) { // condicional
        http_response_code(429);
        header('Retry-After: 900'); // 15 min
        exit('Demasiados intentos. Vuelva en 15 minutos.'); // termina ejecucion
    }

    $pdo = obtenerPDO();
    $stmt = $pdo->prepare( // prepara consulta SQL (defensa anti-SQLi)
        // hashea contrasena con BCrypt
        // hashea contrasena con BCrypt
        'SELECT id, password_hash, rol, activo, intentos_fallidos
         FROM usuarios WHERE email = :em LIMIT 1');
    $stmt->execute([':em' => $email]); // ejecuta la consulta con parametros enlazados
    $u = $stmt->fetch(PDO::FETCH_ASSOC); // recupera la siguiente fila como array

    // password_verify usa tiempo constante (resistente a timing attacks)
    // hashea contrasena con BCrypt
    // hashea contrasena con BCrypt
    if (!$u || !$u['activo'] || !password_verify($pass, $u['password_hash'])) {
        $_SESSION[$intentosKey] = ($_SESSION[$intentosKey] ?? 0) + 1;
        http_response_code(401);
        echo 'Credenciales invalidas'; // imprime al output del servidor
        exit;
    }

    // CRITICAL: regenerar el ID para prevenir session fixation
    session_regenerate_id(true); // cambia el ID para evitar session fixation

    $_SESSION['user_id']    = (int)$u['id'];
    $_SESSION['user_rol']   = $u['rol'];
    $_SESSION['login_time'] = time();
    $_SESSION['ip']         = $_SERVER['REMOTE_ADDR'];
    $_SESSION['ua']         = $_SERVER['HTTP_USER_AGENT'] ?? '';
    unset($_SESSION[$intentosKey]);

    // Auditoria
    error_log("[LOGIN] user_id={$u['id']} ip={$_SERVER['REMOTE_ADDR']}");

    header('Location: /dashboard.php'); // envia cabecera HTTP al cliente
    exit;
}

// Generar token CSRF para mostrar el formulario
if (empty($_SESSION['_csrf'])) { // condicional
    $_SESSION['_csrf'] = bin2hex(random_bytes(32));
}
$csrf = $_SESSION['_csrf'];
?>
<!DOCTYPE html>
<html lang="es-EC">
<head><meta charset="UTF-8"><title>Login UTEQ</title></head>
<body>
<form method="POST">
  // escapa caracteres HTML (defensa anti-XSS)
  // escapa caracteres HTML (defensa anti-XSS)
  <input type="hidden" name="_csrf" value="<?= htmlspecialchars($csrf) ?>">
  <label>Email: <input type="email" name="email" required></label>
  <label>Password: <input type="password" name="password" required minlength="8"></label>
  <button type="submit">Ingresar</button>
</form>
</body></html>
\end{lstlisting}

\textbf{Defensas implementadas en este código (todas combinadas):}
\begin{enumerate}
\item Cookies con \texttt{Secure}, \texttt{HttpOnly},
\texttt{SameSite=Strict}.
\item Token CSRF generado por sesión y validado con
\texttt{hash\_equals} (tiempo constante).
\item Validación email con \texttt{FILTER\_VALIDATE\_EMAIL}.
\item Rate limiting (5 intentos por IP).
\item \texttt{password\_verify} (BCrypt, resistente a timing attacks).
\item \texttt{session\_regenerate\_id(true)}: previene session
fixation.
\item Auditoría con \texttt{error\_log}.
\item \texttt{htmlspecialchars} al renderizar (anti-XSS).
\item Verificación de cuenta activa.
\end{enumerate}
\end{cajaejemplo}

\section{Almacenamiento del cliente: localStorage, sessionStorage e IndexedDB}

Además del almacenamiento del servidor (cookies de sesión,
JWT), el navegador ofrece tres APIs nativas para guardar datos del
lado cliente: \texttt{localStorage}, \texttt{sessionStorage} e
\texttt{IndexedDB}. Conocer cuándo usar cada una es esencial para
SPAs modernas.


\subsection{Comparativa de las APIs de almacenamiento del cliente}

La tabla siguiente contrasta las tres APIs de almacenamiento
cliente por capacidad, persistencia, tipo de datos y casos de uso
recomendados.


\begin{table}[htbp]
\centering
\caption{Almacenamiento del cliente en navegadores modernos.}
\footnotesize
\begin{tabularx}{\textwidth}{lXll}
\toprule
\textbf{Mecanismo} & \textbf{Características} & \textbf{Capacidad} & \textbf{API}\\
\midrule
Cookies & Enviadas con cada petición; \emph{HttpOnly} oculta de JS & 4 KB & \texttt{document.cookie}\\
\texttt{localStorage} & Persistente, mismo origen, síncrono, solo strings & 5--10 MB & Síncrona\\
\texttt{sessionStorage} & Solo durante la pestaña, mismo origen & 5--10 MB & Síncrona\\
IndexedDB & Persistente, transaccional, asíncrono, objetos & GB & Asíncrona\\
Cache API & Almacena objetos \texttt{Response} (Service Workers) & GB & Asíncrona\\
\bottomrule
\end{tabularx}
\end{table}

\begin{cajaadvertencia}{La regla de oro --- ¿dónde guardar el JWT?}
\textbf{Nunca} guarde el JWT en \texttt{localStorage} si su aplicación
tiene cualquier riesgo de XSS (lo cual es prácticamente toda
aplicación). El JavaScript puede leer \texttt{localStorage}; un solo
script malicioso inyectado roba todos los tokens.

\textbf{Solución correcta:} cookie \texttt{HttpOnly} + \texttt{Secure}
+ \texttt{SameSite=Strict}. JavaScript no puede leer cookies
\texttt{HttpOnly}; el navegador las envía automáticamente al servidor
en cada petición al mismo origen.
\end{cajaadvertencia}

\subsection{Ejemplos prácticos de almacenamiento}

Los siguientes ejemplos breves muestran las tres APIs en
acción para el caso más frecuente: persistir preferencias de UI.

\textbf{Ejemplo corto comentado (localStorage) (Listado~\ref{lst:10.13}):}
\begin{lstlisting}[style=js,numbers=none,caption={Ejemplo de JavaScript},label=lst:10.13]
// localStorage persiste INDEFINIDAMENTE (hasta que el usuario lo borre)
// Solo guarda STRINGS, asi que objetos hay que serializarlos con JSON
localStorage.setItem('tema', 'oscuro'); // escribe
const tema = localStorage.getItem('tema'); // lee
localStorage.removeItem('tema'); // borra una clave
localStorage.clear(); // borra todo

// Guardar un objeto: serializar con JSON.stringify
const prefs = { tema: 'oscuro', idioma: 'es' }; // constante (no reasignable)
// serializa objeto a string JSON
// serializa objeto a string JSON
localStorage.setItem('preferencias', JSON.stringify(prefs));
// Recuperar y parsear
// constante (no reasignable)
// constante (no reasignable)
const guardado = JSON.parse(localStorage.getItem('preferencias') || '{}');
\end{lstlisting}


\begin{cajaejemplo}{Ejemplo 10.6 --- Persistencia de preferencias UI con localStorage}

\textbf{Enunciado.} Persistir preferencias de UI (tema oscuro, idioma) en \texttt{localStorage} desde JavaScript, mostrando cómo serializar objetos con JSON y manejar errores de cuota.

\begin{lstlisting}[style=js,caption={preferencias.js},label=lst:10.14]
"use strict";

const PREF_KEY = "uteq:preferencias:v1"; // constante (no reasignable)

const Preferencias = { // constante (no reasignable)
  cargar() {
    try {
      const raw = localStorage.getItem(PREF_KEY); // constante (no reasignable)
      return raw ? JSON.parse(raw) : this.defectos(); // devuelve valor de la funcion
    } catch (e) {
      console.warn("Error leyendo preferencias", e);
      return this.defectos(); // devuelve valor de la funcion
    }
  },
  guardar(p) {
    try {
      localStorage.setItem(PREF_KEY, JSON.stringify(p)); // serializa objeto a string JSON
    } catch (e) {
      console.error("Error guardando preferencias", e);
    }
  },
  defectos() {
    return { tema: "auto", idioma: "es-EC", // devuelve valor de la funcion
             fontSize: 1.0, animaciones: true };
  },
  aplicar(p) {
    document.documentElement.dataset.tema = p.tema;
    document.documentElement.lang = p.idioma;
    document.documentElement.style.fontSize = `${p.fontSize}rem`;
    if (!p.animaciones) { // condicional
      document.documentElement.classList.add("sin-animaciones");
    }
  }
};

// Al cargar
document.addEventListener("DOMContentLoaded", () => { // registra manejador de evento
  const p = Preferencias.cargar(); // constante (no reasignable)
  Preferencias.aplicar(p);

  // Boton de cambio de tema
  // busca elemento por id
  // busca elemento por id
  document.getElementById("toggle-tema")?.addEventListener("click", () => {
    p.tema = p.tema === "claro" ? "oscuro" : "claro";
    Preferencias.guardar(p);
    Preferencias.aplicar(p);
  });
});
\end{lstlisting}
\end{cajaejemplo}

\section{Práctica integrada: carrito de compras con sesiones}

La práctica de la semana es construir un carrito de compras
completo con sesiones de servidor, integrando todos los conceptos:
Request, Response, Session, Cookie segura y validación.


\begin{cajaejercicio}{Ejercicio 10.1 --- Carrito de compras stateful en las tres plataformas (Sumativa GA-Práctica)}

\textbf{Objetivo:} implementar un carrito de compras con sesiones de
servidor en PHP 8.4, ASP.NET Core 10 y Spring Boot 4, evidenciando las
diferencias de API y cómo cada plataforma maneja seguridad de cookies.

\textbf{Especificación funcional:}
\begin{itemize}
\item Página de listado con productos hardcodeados (nombre, precio).
\item Botón «Agregar» para cada producto que lo añade al carrito.
\item Página del carrito que muestra los productos agregados y total.
\item Botón «Eliminar» por producto en el carrito.
\item Botón «Vaciar carrito».
\item La cookie de sesión debe tener \texttt{HttpOnly},
\texttt{Secure} y \texttt{SameSite=Strict}.
\end{itemize}

\subsubsection*{IDE recomendado}
\textbf{IntelliJ IDEA Ultimate} con plugins:
\begin{itemize}
\item \emph{PHP} (JetBrains, instalado por defecto en Ultimate).
\item \emph{Spring Boot} (incluido en Ultimate).
\item \emph{.NET Core} (Marketplace) o usar JetBrains Rider para C\#.
\end{itemize}

\subsubsection*{Plataforma 1 --- PHP 8.4 con XAMPP}

\textbf{Paso 1:} verificar XAMPP. Iniciar el panel
(\texttt{xampp-control}). Pulsar \texttt{Start} en Apache. Verificar
en \url{http://localhost} que aparece el dashboard XAMPP.

\textbf{Paso 2:} crear carpeta \texttt{C:\textbackslash{}xampp\textbackslash{}htdocs\textbackslash{}carrito-php}.
En IntelliJ: \texttt{File $\rightarrow$ Open}, seleccionar la carpeta.

\textbf{Paso 3:} crear \texttt{config.php} (Listado~\ref{lst:10.15}):

\begin{lstlisting}[style=php,caption={config.php},label=lst:10.15]
<?php
declare(strict_types=1);

ini_set('session.cookie_httponly', '1');
ini_set('session.cookie_secure',   '1');
ini_set('session.cookie_samesite', 'Strict');
ini_set('session.use_strict_mode', '1');
ini_set('session.gc_maxlifetime',  '1800');

session_name('CARRITO_PHP');
session_start(); // arranca o reanuda la sesion

// Inicializar carrito si no existe
if (!isset($_SESSION['carrito'])) { // condicional
    $_SESSION['carrito'] = [];
}

// Catalogo hardcodeado
const CATALOGO = [
    1 => ['id'=>1, 'nombre'=>'Cafe galletti 500g', 'precio'=>4.50],
    2 => ['id'=>2, 'nombre'=>'Cacao Pacari 70%',   'precio'=>6.00],
    3 => ['id'=>3, 'nombre'=>'Te verde organico',  'precio'=>3.20],
    4 => ['id'=>4, 'nombre'=>'Panela en bloque',   'precio'=>2.80],
];
\end{lstlisting}

\textbf{Paso 4:} crear \texttt{index.php} (Listado~\ref{lst:10.16}):

\begin{lstlisting}[style=php,caption={index.php},label=lst:10.16]
<?php require_once 'config.php'; ?>
<!DOCTYPE html>
<html lang="es-EC">
<head><meta charset="UTF-8"><title>Tienda UTEQ</title>
<style>
  body{font-family:system-ui;max-width:700px;margin:2rem auto;padding:0 1rem}
  h1{color:#006633} table{width:100%;border-collapse:collapse}
  td,th{padding:.5rem;text-align:left;border-bottom:1px solid #eee}
  .btn{padding:.4rem .8rem;background:#006633;color:white;
       border:0;border-radius:4px;cursor:pointer}
</style>
</head>
<body>
<h1>Productos UTEQ</h1>
<p><a href="carrito.php">Ver carrito (<?= count($_SESSION['carrito']) ?>)</a></p>
<table>
  <thead><tr><th>Producto</th><th>Precio</th><th></th></tr></thead>
  <tbody>
  <?php foreach (CATALOGO as $p): ?>
    <tr>
      <td><?= htmlspecialchars($p['nombre']) ?></td> // escapa caracteres HTML (defensa anti-XSS)
      <td>$<?= number_format($p['precio'], 2) ?></td>
      <td>
        <form action="agregar.php" method="post" style="margin:0">
          <input type="hidden" name="id" value="<?= $p['id'] ?>">
          <button type="submit" class="btn">Agregar</button>
        </form>
      </td>
    </tr>
  <?php endforeach; ?>
  </tbody>
</table>
</body></html>
\end{lstlisting}

\textbf{Paso 5:} crear \texttt{agregar.php}, \texttt{eliminar.php},
\texttt{vaciar.php}, \texttt{carrito.php} (continuación obvia siguiendo
el patrón).

\textbf{Paso 6:} probar.
\begin{enumerate}
\item Abrir \url{http://localhost/carrito-php/index.php}.
\item Agregar 3 productos diferentes.
\item Ir a Carrito: aparecen los 3.
\item Cerrar la pestaña, abrir nueva: el carrito persiste.
\item Cerrar el navegador completamente, abrirlo de nuevo: el carrito
sigue ahí (porque \texttt{lifetime=1800}).
\end{enumerate}

\textbf{Paso 7:} verificar la cookie. F12 $\rightarrow$ Application
$\rightarrow$ Cookies $\rightarrow$ \url{http://localhost}. Debe
aparecer \texttt{CARRITO\_PHP} con: \texttt{HttpOnly = true},
\texttt{Secure = true}, \texttt{SameSite = Strict}.

\textbf{Resultado esperado:} carrito persistente entre páginas y entre
sesiones del navegador, con cookie hardenizada visible en DevTools.

\subsubsection*{Plataforma 2 --- ASP.NET Core 10}

\textbf{Paso 1:} verificar .NET 8. \texttt{dotnet --version} debe
devolver \texttt{8.0.x}.

\textbf{Paso 2:} crear el proyecto desde la terminal de IntelliJ
(\texttt{Alt+F12}) (Listado~\ref{lst:10.17}):
\begin{lstlisting}[style=shell,numbers=none,caption={Ejemplo de Shell},label=lst:10.17]
dotnet new mvc -n CarritoAsp -o CarritoAsp
cd CarritoAsp # cambia directorio
\end{lstlisting}

\textbf{Paso 3:} configurar sesiones en \texttt{Program.cs} (Listado~\ref{lst:10.18}):

\begin{lstlisting}[style=cs,caption={Program.cs},label=lst:10.18]
using Microsoft.AspNetCore.Builder; // importa namespace
using Microsoft.Extensions.DependencyInjection; // importa namespace

var builder = WebApplication.CreateBuilder(args); // crea el builder de la app

builder.Services.AddControllersWithViews(); // configuracion de servicios DI
builder.Services.AddDistributedMemoryCache(); // configuracion de servicios DI
builder.Services.AddSession(o => { // configuracion de servicios DI
    o.IdleTimeout = TimeSpan.FromMinutes(30);
    o.Cookie.Name        = "CARRITO_ASP";
    o.Cookie.HttpOnly    = true;
    o.Cookie.IsEssential = true;
    o.Cookie.SecurePolicy = CookieSecurePolicy.Always;
    o.Cookie.SameSite     = SameSiteMode.Strict;
});

var app = builder.Build(); // construye la app a partir del builder
app.UseHttpsRedirection();
app.UseStaticFiles();
app.UseRouting(); // activa el routing del pipeline
app.UseSession();
app.UseAuthorization();
app.MapControllerRoute(
    name: "default",
    pattern: "{controller=Tienda}/{action=Index}/{id?}");
app.Run(); // arranca la aplicacion
\end{lstlisting}

\textbf{Paso 4:} crear \texttt{Models/Producto.cs} y
\texttt{Controllers/TiendaController.cs} (con acciones
\texttt{Index}, \texttt{Agregar}, \texttt{Carrito},
\texttt{Eliminar}, \texttt{Vaciar} usando
\texttt{HttpContext.Session.SetString} con JSON).

\textbf{Paso 5:} ejecutar con \texttt{dotnet run}. Abrir
\url{https://localhost:5001}.

\textbf{Paso 6:} verificar la cookie \texttt{CARRITO\_ASP} en
DevTools: debe tener \texttt{HttpOnly}, \texttt{Secure},
\texttt{SameSite=Strict}.

\subsubsection*{Plataforma 3 --- Spring Boot 4}

\textbf{Paso 1:} en IntelliJ Ultimate:
\texttt{File $\rightarrow$ New $\rightarrow$ Project $\rightarrow$ Spring Boot}.
Java 21, Spring Boot 4, dependencias \texttt{Spring Web} y
\texttt{Thymeleaf}.

\textbf{Paso 2:} configurar \texttt{src/main/resources/application.yml} (Listado~\ref{lst:10.19}):

\begin{lstlisting}[style=yaml,caption={application.yml},label=lst:10.19]
server: # configuracion del servidor embebido
  port: 8080 # puerto donde escucha la aplicacion
  servlet:
    session:
      timeout: 30m
      cookie:
        name:      CARRITO_SPRING
        http-only: true
        secure:    false # true en HTTPS de produccion
        same-site: strict
\end{lstlisting}

\textbf{Paso 3:} crear \texttt{Producto.java} (record \texttt{Serializable}),
\texttt{TiendaController.java} con \texttt{HttpSession}.

\textbf{Paso 4:} ejecutar y probar en
\url{http://localhost:8080}.

\subsubsection*{Tabla comparativa de la práctica (entrega)}

Cada estudiante debe completar y entregar:

\begin{tabularx}{\textwidth}{lXXX}
\toprule
\textbf{Aspecto} & \textbf{PHP} & \textbf{ASP.NET} & \textbf{Spring}\\
\midrule
Líneas de código del proyecto & & & \\
Tiempo de arranque & & & \\
Cookie tiene HttpOnly & ¿sí/no? & ¿sí/no? & ¿sí/no?\\
Cookie tiene Secure & & & \\
Cookie tiene SameSite=Strict & & & \\
Carrito persiste entre pestañas & & & \\
\bottomrule
\end{tabularx}

\subsubsection*{Reflexión final}

Responder por escrito en máximo 200 palabras: ¿qué plataforma le
resultó más fácil? ¿Cuál considera más segura por defecto? ¿Cuál usaría
en un proyecto real y por qué?

\textbf{Esta es la \emph{Sumativa GA-Práctica en clase} de la semana
10.}
\end{cajaejercicio}

\section{Buenas prácticas profesionales en gestión de estado}

Las técnicas de gestión de estado tienen su propio conjunto
de reglas profesionales. La caja siguiente las reúne en 15 puntos.


\begin{cajabuenas}{Quince reglas profesionales 2026}
\begin{enumerate}
\item Cookies de sesión SIEMPRE con \texttt{Secure},
\texttt{HttpOnly}, \texttt{SameSite=Strict}.
\item Regenerar el ID de sesión al hacer login (anti-fixation).
\item Tiempo de inactividad razonable: 15--30 min para apps
sensibles, 60 min para uso general.
\item Logout debe invalidar el ID de sesión, no solo limpiar variables.
\item Para APIs REST que sirven SPAs: preferir JWT en cookie
\texttt{HttpOnly}.
\item NUNCA guardar JWT en \texttt{localStorage} en sitios con riesgo
de XSS.
\item Validar token CSRF en todo POST/PUT/PATCH/DELETE.
\item En clusters: usar Redis o BD compartida para sesiones, no
memoria local.
\item No almacenar información sensible en la cookie misma; solo el ID.
\item Auditar logins exitosos y fallidos con IP y User-Agent.
\item Implementar rate limiting (máx. 5--10 intentos por IP por minuto).
\item Cookies con \texttt{Domain} restrictivo (no
\texttt{example.com} si solo se necesita \texttt{app.example.com}).
\item Nombres de cookies que NO revelen la tecnología
(\texttt{SESS} en lugar de \texttt{PHPSESSID}).
\item En SPAs Angular/React: usar \texttt{credentials: 'include'} y
configurar CORS con \texttt{Access-Control-Allow-Credentials: true}.
\item Documentar la estrategia de gestión de estado en la ADR-001 del
proyecto.
\end{enumerate}
\end{cajabuenas}

\section{Ejercicio resuelto adicional con resultado visual}

Como cierre de la semana, el siguiente ejercicio adicional
está completamente resuelto e incluye captura del resultado visual.
Se complementa con un ejercicio propuesto para que el estudiante
practique.


\begin{cajaejemplo}{Ejercicio resuelto 10.1 --- Contador de visitas con sesión PHP}
\textbf{Enunciado.} Implementar una página PHP que muestre cuántas
veces el usuario ha visitado el sitio en la sesión actual, con
botones para incrementar manualmente y para reiniciar el contador.
La cookie de sesión debe llevar todos los flags de seguridad de
producción.

\textbf{Código completo (\texttt{contador.php}):}

\begin{lstlisting}[language=PHP,caption={Contador stateful con sesión segura},label=lst:10.20]
<?php
// Configuracion ANTES de session_start()
session_set_cookie_params([ // flags de cookie
  'lifetime' => 1800,
  'path'     => '/',
  'domain'   => '',
  'secure'   => false, // true en HTTPS de produccion
  'httponly' => true,
  'samesite' => 'Strict'
]);
session_name('VISITAS_SESS');
session_start(); // arranca o reanuda la sesion

// Inicializar contador
if (!isset($_SESSION['visitas'])) { // condicional
    $_SESSION['visitas'] = 1;
} else {
    // Solo incrementar en GET sin parametros (visita real)
    // condicional
    // condicional
    if ($_SERVER['REQUEST_METHOD'] === 'GET' && empty($_GET)) {
        $_SESSION['visitas']++;
    }
}

// Acciones
if (isset($_GET['accion'])) { // condicional
    // condicional
    // condicional
    if ($_GET['accion'] === 'incrementar') $_SESSION['visitas']++;
    // condicional
    // condicional
    if ($_GET['accion'] === 'reiniciar')   $_SESSION['visitas'] = 0;
    header('Location: contador.php'); // envia cabecera HTTP al cliente
    exit;
}
?>
<!DOCTYPE html>
<html lang="es">
<head><meta charset="UTF-8"><title>Contador de visitas</title></head>
<body>
  <h1>Contador de visitas</h1>
  <p>Visitas en esta sesion: <strong><?= $_SESSION['visitas'] ?></strong></p>
  // escapa caracteres HTML (defensa anti-XSS)
  // escapa caracteres HTML (defensa anti-XSS)
  <p>Session ID: <code><?= htmlspecialchars(session_id()) ?></code></p>
  <a href="?accion=incrementar">Incrementar +1</a> |
  <a href="?accion=reiniciar">Reiniciar</a>
</body>
</html>
\end{lstlisting}

\textbf{Resultado visual en el navegador} (simulación de cómo se
verá la página en la pantalla del estudiante tras varias recargas
e incrementos manuales):

\begin{center}
\fbox{\begin{minipage}{0.85\textwidth}
\vspace{0.2cm}
\hspace{0.3cm}\textbf{\large Contador de visitas}\\[0.4cm]
\hspace{0.3cm}Visitas en esta sesi\'on: \textbf{7}\\[0.2cm]
\hspace{0.3cm}Session ID: \texttt{a8f3d9c2e1b4...}\\[0.3cm]
\hspace{0.3cm}\textcolor{blue}{\underline{Incrementar +1}} \textbar{}
\textcolor{blue}{\underline{Reiniciar}}\\
\vspace{0.2cm}
\end{minipage}}
\end{center}

\textbf{Verificación de los flags de seguridad} con DevTools del
navegador (pestaña \emph{Application $\to$ Cookies}):

\begin{center}
\fbox{\begin{minipage}{0.85\textwidth}
\vspace{0.15cm}\footnotesize
\hspace{0.3cm}\textbf{Cookie:} \texttt{VISITAS\_SESS}\\
\hspace{0.3cm}Value: \texttt{a8f3d9c2e1b4...}\\
\hspace{0.3cm}HttpOnly: $\checkmark$ \quad Secure: $\times$ \quad SameSite: \texttt{Strict}\\
\hspace{0.3cm}Expires: Session\\
\vspace{0.15cm}
\end{minipage}}
\end{center}

\textbf{Discusión.} El contador demuestra que la sesión persiste
entre recargas (el navegador reenvía la cookie automáticamente) y
que el servidor reconoce al mismo usuario. La cookie es invisible a
JavaScript (\texttt{HttpOnly}), bloqueando ataques XSS de robo de
sesión, y solo se envía desde el propio dominio
(\texttt{SameSite=Strict}), bloqueando CSRF. En producción HTTPS,
añadir \texttt{secure=true} bloquea su envío por HTTP plano.
\end{cajaejemplo}

\begin{cajaejercicio}{Ejercicio propuesto 10.2 --- Autenticación con login persistente «Recuérdame»}
\textbf{Enunciado.} Implementar en cualquiera de las tres plataformas
(PHP, Spring Boot o ASP.NET Core) un flujo de login con la opción
\emph{«Recuérdame»} que mantenga al usuario autenticado durante 30
días si la marca, o solo durante la sesión del navegador si no.

\textbf{Requisitos técnicos:}
\begin{enumerate}
\item Formulario de login con campos \texttt{usuario},
\texttt{contraseña} y checkbox \texttt{recordarme}.
\item Si \texttt{recordarme} está marcado: crear una cookie
adicional \texttt{REMEMBER\_TOKEN} con un valor aleatorio de 256 bits
(\texttt{random\_bytes(32)} en PHP, \texttt{SecureRandom} en Java),
\texttt{Max-Age=2592000} (30 días), todos los flags de seguridad.
\item Persistir el hash del token (\texttt{password\_hash()} o
\texttt{BCrypt}) en una tabla \texttt{remember\_tokens(user\_id,
token\_hash, expires\_at)}.
\item En cada nueva sesión, si el usuario no está autenticado pero
existe \texttt{REMEMBER\_TOKEN}, validar contra la tabla y
autenticar automáticamente.
\item Al hacer logout, eliminar la fila correspondiente y borrar la
cookie con \texttt{Max-Age=0}.
\item Implementar rotación del token tras cada uso (defensa contra
robo).
\end{enumerate}

\textbf{Criterios de evaluación:}
\begin{itemize}
\item Token nunca se almacena en plano (solo el hash).
\item Cookies con todos los flags (\texttt{HttpOnly},
\texttt{Secure}, \texttt{SameSite=Strict}, \texttt{Max-Age}).
\item Rotación correcta del token tras autenticación automática.
\item Caso de logout limpia BD y navegador.
\item Documentación con capturas de Postman o navegador mostrando
las cookies generadas.
\end{itemize}

\textbf{Lecturas recomendadas:} \cite{jones2015jwt} para
comparar con JWT; \cite{ortega2022securecoding} para defensas
contra robo de sesión.
\end{cajaejercicio}

% =====================================================================
%       SEMANA 11: ACCESO A DATOS DESDE APLICACIONES WEB
% =====================================================================
\chapter{Acceso a datos desde aplicaciones web: SQL parametrizado y ORM}
\label{cap:semana11}

\epigraph{«Toda aplicación que persiste datos hereda, tarde o
temprano, la disciplina o el caos de cómo accede a ellos. Elegir bien
es proteger el sistema durante años.»}{--- \textit{Adaptado de Martin Fowler}}

\section*{Resultado de aprendizaje de la semana}
Al concluir la semana 11, el estudiante diseña e implementa el acceso a
datos relacionales desde aplicaciones web utilizando SQL parametrizado
y mapeo objeto-relacional (ORM), aplicando los patrones Repository y
Unit of Work, prevenir totalmente la inyección SQL y comprende el
problema N+1 y su solución.

\section{Historia y panorama del acceso a datos}

El acceso a datos es probablemente la decisión arquitectónica
con mayor impacto en el rendimiento y la mantenibilidad de una
aplicación web. La elección entre SQL plano, JPA/Hibernate, Entity
Framework Core o Eloquent no es de detalle: condiciona todo el código
del backend.


\begin{cajahistoria}{Del SQL embebido a los ORM modernos --- 50 años de iteración}
La historia del acceso a datos desde aplicaciones web comienza con SQL
en su forma más cruda: cadenas de texto concatenadas dentro del código
de la aplicación. En 1995--2005 era habitual ver código PHP así:
\texttt{\$sql = \char34 SELECT * FROM users WHERE name = '\char34{} . \$\_POST['name'] . \char34 '\char34;}.
Este patrón --- inocente en apariencia --- destruyó la seguridad de
miles de aplicaciones durante una década por SQL Injection.

La primera respuesta fueron las \textbf{sentencias preparadas}
(\emph{prepared statements}), formalizadas en SQL-92 y ampliamente
disponibles desde principios de los 2000: el SQL se compila una vez
con marcadores (\texttt{?} o \texttt{:nombre}), y los valores viajan
separadamente sin interpretarse como código. Esto cerró la mayor
puerta de entrada a los ataques.

Una segunda evolución fueron los \textbf{Query Builders}: APIs
fluidas para construir SQL programáticamente sin escribirlo a mano
(\texttt{Doctrine DBAL} en PHP, \texttt{jOOQ} en Java). Permiten
componer consultas dinámicas sin riesgo de inyección.

El salto cualitativo vino con los \textbf{ORM} (Object-Relational
Mappers), que abstraen completamente la base de datos detrás de
objetos del lenguaje. Hibernate en Java (Gavin King, 2001), Active
Record en Ruby on Rails (David Heinemeier Hansson, 2004), Doctrine en
PHP (Roman Borschel, 2006), Entity Framework en .NET (Microsoft, 2008)
y Eloquent en Laravel (Taylor Otwell, 2011) representan distintas
filosofías sobre el mismo problema. Hoy un desarrollador web rara vez
escribe SQL a mano; pero comprender lo que hace el ORM por debajo es
indispensable.

La literatura técnica confirma que la elección del mecanismo de
acceso a datos es una decisión arquitectónica de primer orden.
\cite{fowler2002patterns} describe formalmente los patrones que hoy
dominan (Data Mapper, Active Record, Unit of Work, Identity Map);
\cite{ottinger2017hibernate} y \cite{bauer2015java} son las
referencias canónicas para Hibernate; \cite{smith2018entity} para
Entity Framework Core; \cite{otwell2024laravel} para Eloquent. El
trabajo empírico de \cite{salah2017orm} compara performance entre
los principales ORM y documenta que la diferencia de rendimiento
respecto a SQL plano oscila entre 10\% y 30\%, justificándose
ampliamente por la productividad ganada. Estudios más recientes
\cite{ottinger2017hibernate} en revistas indexadas confirman que el problema
N+1 sigue siendo la causa número uno de degradación de performance
en aplicaciones Java EE y Spring Boot en producción.
\end{cajahistoria}

\subsection{Las cuatro estrategias de acceso a datos}

En la industria conviven cuatro estrategias para acceder a
datos desde una aplicación web. La tabla siguiente las contrasta
con sus ventajas, desventajas y cuándo elegir cada una.


\begin{table}[htbp]
\centering
\caption{Estrategias de acceso a datos en aplicaciones web (2026).}
\footnotesize
\begin{tabularx}{\textwidth}{lXX}
\toprule
\textbf{Estrategia} & \textbf{Características} & \textbf{Cuándo usarla}\\
\midrule
Driver nativo + SQL & Máximo control y rendimiento; máxima responsabilidad. PDO (PHP), JDBC (Java), ADO.NET (C\#). & Reportes complejos, consultas analíticas, ETL.\\
Query Builder & SQL programático sin concatenación. Doctrine DBAL, jOOQ, knex.js. & Consultas dinámicas con muchos filtros opcionales.\\
ORM Active Record & La entidad \emph{es} su unidad de persistencia: \texttt{User::find(1)->save()}. Eloquent, Rails AR. & CRUDs estándar; prototipado rápido.\\
ORM Data Mapper & Entidades de dominio agnósticas; mapper externo. Hibernate, Doctrine ORM, EF Core. & Aplicaciones empresariales medianas/grandes con DDD.\\
\bottomrule
\end{tabularx}
\end{table}

\subsection{Tabla evolutiva de los ORM principales}

Los ORM dominantes han evolucionado durante dos décadas
hasta sus versiones actuales. La tabla siguiente recoge las
versiones de referencia 2026 con su licencia y características
notables.


\begin{table}[htbp]
\centering
\caption{Versiones actuales de ORM en 2026.}
\footnotesize
\begin{tabular}{llll}
\toprule
\textbf{ORM} & \textbf{Lenguaje} & \textbf{Versión 2026} & \textbf{Filosofía}\\
\midrule
Hibernate / JPA      & Java       & Hibernate 6.6 / JPA 3.2 & Data Mapper\\
Spring Data JPA      & Java       & 3.4                     & Repository sobre JPA\\
Eloquent             & PHP        & Laravel 11 incluido     & Active Record\\
Doctrine ORM         & PHP        & 3.x                     & Data Mapper\\
Entity Framework Core & C\#       & 8.0                     & Code-First, Data Mapper\\
TypeORM              & TypeScript & 0.3                     & Híbrido AR/DM\\
Prisma               & TypeScript & 5.x                     & Schema-first generador\\
SQLAlchemy           & Python     & 2.x                     & Data Mapper + AR\\
Active Record (Rails) & Ruby      & 8.x                     & Active Record canónico\\
GORM                 & Go         & 1.25                    & Active Record\\
\bottomrule
\end{tabular}
\end{table}

\section{Consultas SQL desde aplicaciones web}

Aunque el ORM oculte el SQL, el desarrollador web profesional
debe conocer cómo se ejecutan las consultas desde aplicaciones,
qué amenazas existen y qué defensas debe aplicar siempre.


\subsection{La amenaza fundamental: SQL Injection}

SQL Injection no es una vulnerabilidad teórica: ha sido la
causa de ataques contra Sony, LinkedIn, Yahoo y miles de
organizaciones. Sigue siendo la \emph{Top 3} del OWASP en 2021. Su
mecánica es engañosamente simple.

\textbf{Ejemplo corto comentado (anatomía de un SQL Injection) (Listado~\ref{lst:11.1}):}
\begin{lstlisting}[language=PHP,numbers=none,caption={Ejemplo de PHP},label=lst:11.1]
<?php
// Codigo vulnerable: concatena directamente la entrada del usuario
$email = $_POST['email'];
$sql = "SELECT * FROM users WHERE email = '$email'";

// El atacante introduce en el campo email:
//   ' OR '1'='1
// La consulta resultante se convierte en:
//   SELECT * FROM users WHERE email = '' OR '1'='1'
// Que devuelve TODOS los usuarios. El atacante entra sin password.
?>
\end{lstlisting}


\begin{cajaadvertencia}{Anti-patrón --- el código que NUNCA debe escribirse}
\begin{lstlisting}[style=php,numbers=none,caption={Ejemplo de PHP},label=lst:11.2]
// VULNERABLE - PHP
$sql = "SELECT * FROM users WHERE email = '" . $_POST['email'] . "'";
$res = $pdo->query($sql);
\end{lstlisting}
\textbf{Ejemplo de ataque trivial:}\\
Si \texttt{\$\_POST['email'] = \char34{}admin' OR '1'='1' -{}-\char34{}}, la consulta
ejecutada será:\\
\texttt{SELECT * FROM users WHERE email = 'admin' OR '1'='1' -{}-'}\\
que devuelve \emph{todos} los usuarios. Variantes más sofisticadas
permiten dropear tablas, exfiltrar contraseñas, leer archivos del
servidor.

\textbf{Estadística reveladora:} en el \emph{Verizon Data Breach
Investigations Report 2025} \cite{verizon2025dbir}, SQL Injection
sigue siendo responsable del 14\% de las brechas de aplicaciones web,
30 años después de su descubrimiento.
\end{cajaadvertencia}

\subsection{La defensa: sentencias preparadas en las tres plataformas}

La defensa universal contra SQLi es la \emph{sentencia
preparada}: el SQL y los datos viajan por canales separados, de
modo que los datos no pueden alterar la estructura del SQL. Las
tres plataformas del curso la implementan.


\subsubsection{PHP con PDO}

PHP dispone de PDO (\emph{PHP Data Objects}), una capa de
abstracción que soporta múltiples motores (MySQL, PostgreSQL, SQLite,
SQL Server) con una API uniforme.


\begin{cajaejemplo}{Ejemplo 11.1 --- PDO con parámetros nominales y posicionales}

\textbf{Enunciado.} Implementar el CRUD de la entidad \texttt{Libro} con Spring Data JPA y H2: definir entidad, repositorio y controlador REST. Se exponen las cuatro operaciones básicas.

\begin{lstlisting}[style=php,caption={UsuarioRepository.php},label=lst:11.3]
<?php
declare(strict_types=1);

class UsuarioRepository // declara clase PHP
{
    public function __construct(private PDO $pdo) {} // metodo publico

    /** Busqueda con parametros NOMINALES (recomendado). */
    public function buscarPorEmail(string $email): ?array { // metodo publico
        $sql = 'SELECT id, nombre, email, rol, activo
                FROM usuarios
                WHERE email = :em AND activo = 1';
        $stmt = $this->pdo->prepare($sql); // prepara consulta SQL (defensa anti-SQLi)
        $stmt->execute([':em' => $email]); // ejecuta la consulta con parametros enlazados
        $u = $stmt->fetch(PDO::FETCH_ASSOC); // recupera la siguiente fila como array
        return $u ?: null; // retorna valor
    }

    /** Busqueda con parametros POSICIONALES. */
    public function buscarPorId(int $id): ?array { // metodo publico
        $stmt = $this->pdo->prepare( // prepara consulta SQL (defensa anti-SQLi)
            'SELECT id, nombre, email FROM usuarios WHERE id = ?');
        $stmt->execute([$id]); // ejecuta la consulta con parametros enlazados
        return $stmt->fetch(PDO::FETCH_ASSOC) ?: null; // recupera la siguiente fila como array
    }

    /** Busqueda con LIKE seguro - nunca concatenar el % */
    // metodo publico
    // metodo publico
    public function buscarPorNombre(string $patron): array {
        $stmt = $this->pdo->prepare( // prepara consulta SQL (defensa anti-SQLi)
            'SELECT id, nombre, email FROM usuarios
             WHERE nombre LIKE :p
             ORDER BY nombre LIMIT 50');
        // El % se anade aqui, NO concatenado en el SQL
        $stmt->execute([':p' => '%' . $patron . '%']); // ejecuta la consulta con parametros enlazados
        return $stmt->fetchAll(PDO::FETCH_ASSOC); // recupera todas las filas restantes
    }

    /** INSERT con retorno del ID generado. */
    public function crear(string $email, string $nombre, // metodo publico
                          string $passwordPlano): int {
        // hashea contrasena con BCrypt
        // hashea contrasena con BCrypt
        $hash = password_hash($passwordPlano, PASSWORD_BCRYPT,
                              ['cost' => 12]);
        $stmt = $this->pdo->prepare( // prepara consulta SQL (defensa anti-SQLi)
            // hashea contrasena con BCrypt
            // hashea contrasena con BCrypt
            'INSERT INTO usuarios (email, nombre, password_hash, rol)
             VALUES (:em, :nom, :hash, :rol)');
        $stmt->execute([ // ejecuta la consulta con parametros enlazados
            ':em'   => $email,
            ':nom'  => $nombre,
            ':hash' => $hash,
            ':rol'  => 'ROLE_USER',
        ]);
        return (int)$this->pdo->lastInsertId(); // retorna valor
    }

    /** Transaccion explicita para operaciones multi-tabla. */
    // metodo publico
    // metodo publico
    public function transferirPuntos(int $de, int $a, int $monto): bool {
        try {
            $this->pdo->beginTransaction();
            $stmt1 = $this->pdo->prepare( // prepara consulta SQL (defensa anti-SQLi)
                'UPDATE saldos SET puntos = puntos - ? WHERE user_id = ?');
            $stmt1->execute([$monto, $de]); // ejecuta la consulta con parametros enlazados
            $stmt2 = $this->pdo->prepare( // prepara consulta SQL (defensa anti-SQLi)
                'UPDATE saldos SET puntos = puntos + ? WHERE user_id = ?');
            $stmt2->execute([$monto, $a]); // ejecuta la consulta con parametros enlazados
            $this->pdo->commit();
            return true; // retorna valor
        } catch (\Throwable $e) {
            $this->pdo->rollBack();
            error_log("Error transferir: " . $e->getMessage());
            return false; // retorna valor
        }
    }
}
\end{lstlisting}

\textbf{Configuración inicial obligatoria del PDO (Listado~\ref{lst:11.4}):}
\begin{lstlisting}[style=php,numbers=none,caption={Ejemplo de PHP},label=lst:11.4]
$pdo = new PDO( // conecta a la base de datos via PDO
    'mysql:host=localhost;dbname=appweb;charset=utf8mb4',
    'usuario', 'password',
    [
        // Errores como excepciones (NO codigos de error silenciosos)
        PDO::ATTR_ERRMODE            => PDO::ERRMODE_EXCEPTION,
        // FETCH_ASSOC por defecto (arrays con keys, sin duplicar)
        PDO::ATTR_DEFAULT_FETCH_MODE => PDO::FETCH_ASSOC,
        // Prepared statements REALES (no emulados, mas seguros)
        PDO::ATTR_EMULATE_PREPARES   => false,
    ]
);
\end{lstlisting}
\end{cajaejemplo}

\subsubsection{Java con JDBC y try-with-resources}

Java accede a bases de datos a través de JDBC. Desde Java 7,
la sintaxis \emph{try-with-resources} garantiza que las conexiones,
\texttt{Statement}s y \texttt{ResultSet}s se cierren automáticamente,
incluso si hay excepciones.


\begin{cajaejemplo}{Ejemplo 11.2 --- JDBC moderno con try-with-resources}

\textbf{Enunciado.} Mostrar el patrón profesional de acceso a datos en Java con JDBC moderno: \texttt{try-with-resources} que gestiona automáticamente el cierre de \texttt{Connection}, \texttt{Statement} y \texttt{ResultSet}.

\begin{lstlisting}[style=java,caption={UsuarioJdbcRepository.java},label=lst:11.5]
package ec.edu.uteq.appweb.repository; // paquete del archivo

import java.sql.*; // importacion de clase externa
import java.util.*; // importacion de clase externa
import javax.sql.DataSource; // importacion de clase externa

public class UsuarioJdbcRepository { // clase publica

    private final DataSource ds; // inmutable

    // metodo publico
    // metodo publico
    public UsuarioJdbcRepository(DataSource ds) { this.ds = ds; }

    public Optional<Usuario> buscarPorEmail(String email) { // metodo publico
        String sql = """
            SELECT id, nombre, email, rol
            FROM usuarios
            WHERE email = ? AND activo = TRUE
            """;
        try (Connection cn = ds.getConnection();
             PreparedStatement ps = cn.prepareStatement(sql)) {
            ps.setString(1, email);
            try (ResultSet rs = ps.executeQuery()) {
                if (rs.next()) { // condicional
                    return Optional.of(new Usuario( // retorna
                        rs.getLong("id"),
                        rs.getString("nombre"),
                        rs.getString("email"),
                        rs.getString("rol")
                    ));
                }
                return Optional.empty(); // retorna
            }
        } catch (SQLException e) {
            // lanza excepcion explicita
            // lanza excepcion explicita
            throw new RepositoryException("Error buscando por email", e);
        }
    }

    public List<Usuario> buscarPorPatron(String patron) { // metodo publico
        String sql = """
            SELECT id, nombre, email, rol
            FROM usuarios
            WHERE nombre LIKE ?
            ORDER BY nombre LIMIT 50
            """;
        List<Usuario> resultado = new ArrayList<>();
        try (Connection cn = ds.getConnection();
             PreparedStatement ps = cn.prepareStatement(sql)) {
            ps.setString(1, "%" + patron + "%");
            try (ResultSet rs = ps.executeQuery()) {
                while (rs.next()) { // bucle while
                    resultado.add(new Usuario(
                        rs.getLong("id"),
                        rs.getString("nombre"),
                        rs.getString("email"),
                        rs.getString("rol")
                    ));
                }
            }
        } catch (SQLException e) {
            // lanza excepcion explicita
            // lanza excepcion explicita
            throw new RepositoryException("Error buscando por patron", e);
        }
        return resultado; // retorna
    }

    // metodo publico
    // metodo publico
    public long crear(String email, String nombre, String hash) {
        String sql = """
            INSERT INTO usuarios (email, nombre, password_hash, rol)
            VALUES (?, ?, ?, 'ROLE_USER')
            """;
        try (Connection cn = ds.getConnection();
             PreparedStatement ps = cn.prepareStatement(sql,
                                       Statement.RETURN_GENERATED_KEYS)) {
            ps.setString(1, email);
            ps.setString(2, nombre);
            ps.setString(3, hash);
            ps.executeUpdate();
            try (ResultSet keys = ps.getGeneratedKeys()) {
                if (keys.next()) return keys.getLong(1); // condicional
                // lanza excepcion explicita
                // lanza excepcion explicita
                throw new RepositoryException("No se obtuvo ID generado");
            }
        } catch (SQLException e) {
            // lanza excepcion explicita
            // lanza excepcion explicita
            throw new RepositoryException("Error creando usuario", e);
        }
    }
}
\end{lstlisting}
\end{cajaejemplo}

\subsubsection{C\# con ADO.NET}

C\# accede a bases de datos a través de ADO.NET, la
biblioteca clásica de .NET. Aunque Entity Framework Core es lo
preferido en proyectos nuevos, ADO.NET sigue siendo la base sobre
la que se construye y aparece en código heredado.


\begin{cajaejemplo}{Ejemplo 11.3 --- ADO.NET puro con SqlClient}

\textbf{Enunciado.} Mostrar el acceso a datos en C\# con ADO.NET puro (\texttt{SqlClient}), antes de pasar a Entity Framework Core. Útil para entender lo que el ORM oculta.

\begin{lstlisting}[style=cs,caption={UsuarioRepository.cs},label=lst:11.6]
using Microsoft.Data.SqlClient; // importa namespace
using AppwebApi.Models; // importa namespace

public class UsuarioRepository // clase publica
{
    private readonly string _connStr; // metodo o miembro privado

    public UsuarioRepository(IConfiguration cfg) => // metodo o miembro publico
        _connStr = cfg.GetConnectionString("Default")!;

    // metodo o miembro publico
    // metodo o miembro publico
    public async Task<Usuario?> BuscarPorEmailAsync(string email)
    {
        const string sql = @"
            SELECT id, nombre, email, rol
            FROM usuarios
            WHERE email = @email AND activo = 1";

        await using var cn = new SqlConnection(_connStr);
        await cn.OpenAsync();

        await using var cmd = new SqlCommand(sql, cn);
        cmd.Parameters.AddWithValue("@email", email);

        await using var rdr = await cmd.ExecuteReaderAsync();
        if (await rdr.ReadAsync()) { // condicional
            return new Usuario { // retorna
                Id     = rdr.GetInt64(0),
                Nombre = rdr.GetString(1),
                Email  = rdr.GetString(2),
                Rol    = rdr.GetString(3)
            };
        }
        return null; // retorna
    }
}
\end{lstlisting}
\end{cajaejemplo}

\section{El patrón Repository: encapsular el acceso a datos}

El patrón Repository encapsula el acceso a datos detrás de
una interfaz limpia, ocultando los detalles de SQL/JPA al resto de
la aplicación. Esto mejora la testabilidad (se puede mockear) y la
mantenibilidad (cambiar de SQL puro a JPA no afecta a los servicios).


\begin{cajadefinicion}{Patrón Repository según Fowler}
\emph{«Mediates between the domain and data mapping layers using a
collection-like interface for accessing domain objects.»}
\cite{fowler2002patterns}

Un Repository expone una API estilo colección sobre las entidades de
dominio, ocultando los detalles del SGBD. Sus responsabilidades:
\begin{itemize}
\item Encapsular consultas (no exponer SQL al servicio).
\item Devolver objetos de dominio o DTOs, no \texttt{ResultSet}.
\item Centralizar la traducción objeto $\leftrightarrow$ tabla.
\item Facilitar el testing con repositorios fake o mocks.
\end{itemize}
\end{cajadefinicion}

\subsection{Repositorio típico en Spring Data JPA}

Spring Data JPA reduce el boilerplate al mínimo: basta con declarar
una interfaz que extiende \texttt{JpaRepository}; Spring genera la
implementación en tiempo de ejecución.

\begin{cajaejemplo}{Ejemplo 11.4 --- ProductoRepository de Spring Data}

\textbf{Enunciado.} Implementar un repositorio Spring Data JPA para la entidad \texttt{Producto} declarando solo la interfaz; Spring genera la implementación con todas las operaciones CRUD.

\begin{lstlisting}[style=java,caption={ProductoRepository.java},label=lst:11.7]
package ec.edu.uteq.appweb.repository; // paquete del archivo

import ec.edu.uteq.appweb.model.Producto; // importacion de clase externa
import org.springframework.data.domain.Page; // importacion de clase externa
import org.springframework.data.domain.Pageable; // importacion de clase externa
import org.springframework.data.jpa.repository.*; // importacion de clase externa
import org.springframework.data.repository.query.Param; // importacion de clase externa
import java.math.BigDecimal; // importacion de clase externa
import java.util.List; // importacion de clase externa

public interface ProductoRepository // declaracion de interfaz
        extends JpaRepository<Producto, Long> {

    // Consulta DERIVADA del nombre del metodo (sin SQL)
    List<Producto> findByStockGreaterThanOrderByNombreAsc(int min);

    Page<Producto> findByCategoriaIgnoreCase(String cat, Pageable pag);

    // JPQL personalizado con parametros nominales
    @Query(""" // JPQL
        SELECT p FROM Producto p
        WHERE p.precio BETWEEN :min AND :max
          AND p.stock > 0
        ORDER BY p.precio ASC
        """)
    List<Producto> rangoPrecio(@Param("min") BigDecimal min,
                               @Param("max") BigDecimal max);

    // Update masivo
    @Modifying // consulta que modifica datos (no SELECT)
    // JPQL
    // JPQL
    @Query("UPDATE Producto p SET p.stock = p.stock - :n WHERE p.id = :id AND p.stock >= :n")
    int decrementarStock(@Param("id") Long id, @Param("n") int n);

    // Native SQL cuando JPQL no alcanza
    @Query(value = """ // JPQL
        SELECT p.*, c.nombre AS cat_nombre
        FROM productos p
        JOIN categorias c ON c.id = p.categoria_id
        WHERE p.stock > 0
        """, nativeQuery = true)
    List<Object[]> reporteStock();
}
\end{lstlisting}
\end{cajaejemplo}

\section{El problema N+1 y su solución}

El problema N+1 es la patología de rendimiento más común en
aplicaciones que usan ORM. Ocurre cuando, al iterar una colección
de N entidades con relaciones, se ejecutan N consultas adicionales
en lugar de una sola con \texttt{JOIN}.


\begin{cajaadvertencia}{Problema N+1 --- el bug silencioso que mata el rendimiento}
Si tiene una entidad \texttt{Producto} con relación
\texttt{@ManyToOne Categoria categoria}, y carga una lista de 100
productos para mostrarlos con su categoría, ocurre lo siguiente sin
optimización:

\begin{enumerate}
\item Una consulta para obtener los 100 productos.
\item Cien consultas adicionales (una por cada producto) para
obtener su categoría asociada.
\end{enumerate}

\textbf{Total: 101 consultas.} Si la página muestra 1000 productos,
son 1001 consultas. La latencia se dispara, la BD colapsa.

\textbf{La solución universal:} \emph{eager loading} explícito, que
hace UN SOLO JOIN en lugar de N+1 consultas.
\end{cajaadvertencia}

\subsection{Solución en cada ORM}

\textbf{Spring Data JPA --- usar JOIN FETCH (Listado~\ref{lst:11.8}):}
\begin{lstlisting}[style=java,numbers=none,caption={Ejemplo de Java},label=lst:11.8]
@Query(""" // JPQL
    SELECT p FROM Producto p
    JOIN FETCH p.categoria
    WHERE p.stock > 0
    """)
List<Producto> findAllConCategoria();
\end{lstlisting}

\textbf{Eloquent (Laravel) --- usar with (Listado~\ref{lst:11.9}):}
\begin{lstlisting}[style=php,numbers=none,caption={Ejemplo de PHP},label=lst:11.9]
$productos = Producto::with('categoria')
    ->where('stock','>',0)
    ->get();
\end{lstlisting}

\textbf{Entity Framework Core --- usar Include (Listado~\ref{lst:11.10}):}
\begin{lstlisting}[style=cs,numbers=none,caption={Ejemplo de C\#},label=lst:11.10]
var productos = await _ctx.Productos
    .Include(p => p.Categoria)
    .Where(p => p.Stock > 0)
    .ToListAsync();
\end{lstlisting}

\textbf{Doctrine ORM --- usar fetch en DQL (Listado~\ref{lst:11.11}):}
\begin{lstlisting}[style=php,numbers=none,caption={Ejemplo de PHP},label=lst:11.11]
$query = $em->createQuery(
    'SELECT p, c FROM App\Entity\Producto p
     JOIN p.categoria c WHERE p.stock > 0');
$productos = $query->getResult();
\end{lstlisting}

\section{Migraciones de esquema con Flyway y Liquibase}

Las migraciones permiten versionar el esquema de la BD como código.
Cada migración es un script SQL numerado que se aplica una sola vez.
La BD «sabe» qué migraciones ha aplicado mediante una tabla de control
(\texttt{flyway\_schema\_history} o \texttt{databasechangelog}).

\begin{cajaejemplo}{Ejemplo 11.5 --- Estructura Flyway en Spring Boot}

\textbf{Enunciado.} Configurar Flyway en un proyecto Spring Boot para gestionar migraciones de base de datos versionadas y reproducibles entre entornos.

\begin{lstlisting}[style=shell,numbers=none,caption={Ejemplo de Shell},label=lst:11.12]
src/main/resources/db/migration/
  V1__crear_usuarios.sql
  V2__crear_productos.sql
  V3__agregar_indice_email_usuarios.sql
  V4__crear_pedidos.sql
\end{lstlisting}

\begin{lstlisting}[style=sql,caption={V1\_\_crear\_usuarios.sql},label=lst:11.13]
-- crea tabla nueva en la BD
CREATE TABLE usuarios (                                     -- crea tabla nueva en la BD
    id            BIGSERIAL PRIMARY KEY,
    -- columna obligatoria (no admite NULL)
    email         VARCHAR(180) NOT NULL UNIQUE,             -- columna obligatoria (no admite NULL)
    -- columna obligatoria (no admite NULL)
    nombre        VARCHAR(100) NOT NULL,                    -- columna obligatoria (no admite NULL)
    -- columna obligatoria (no admite NULL)
    password_hash VARCHAR(255) NOT NULL,                    -- columna obligatoria (no admite NULL)
    -- columna obligatoria (no admite NULL)
    -- columna obligatoria (no admite NULL)
    rol           VARCHAR(30)  NOT NULL DEFAULT 'ROLE_USER',
    -- columna obligatoria (no admite NULL)
    activo        BOOLEAN      NOT NULL DEFAULT TRUE,       -- columna obligatoria (no admite NULL)
    -- columna obligatoria (no admite NULL)
    creado_en     TIMESTAMPTZ  NOT NULL DEFAULT NOW(),      -- columna obligatoria (no admite NULL)
    actualizado_en TIMESTAMPTZ
);

-- crea indice para acelerar busquedas
CREATE INDEX idx_usuarios_email_activo                      -- crea indice para acelerar busquedas
    ON usuarios(email) WHERE activo = TRUE;
\end{lstlisting}

\textbf{Configuración mínima en \texttt{application.yml}:}
\begin{lstlisting}[style=yaml,numbers=none,caption={Ejemplo de YAML},label=lst:11.14]
spring: # configuracion de Spring
  datasource: # configuracion de origen de datos (BD)
    url:      jdbc:postgresql://localhost:5432/appweb
    username: appweb # usuario de la BD
    password: ${DB_PASSWORD} # contrasena (mejor leer de env var)
  jpa: # configuracion JPA/Hibernate
    hibernate: # configuracion especifica de Hibernate
      ddl-auto: validate # NO usar 'update' con Flyway
    properties:
      hibernate.dialect: org.hibernate.dialect.PostgreSQLDialect
  flyway:
    enabled: true
    baseline-on-migrate: true
\end{lstlisting}
\end{cajaejemplo}

\section{Procedimientos almacenados: cuándo aún tiene sentido}

Los procedimientos almacenados (\emph{stored procedures})
fueron durante décadas la forma canónica de encapsular lógica de
datos en el motor de BD. Hoy se usan menos, pero hay casos donde aún
tienen sentido y conviene conocerlos.


\begin{cajacompara}{Procedimientos almacenados vs lógica en aplicación}
\centering\footnotesize
\begin{tabularx}{\textwidth}{lXX}
\toprule
\textbf{Aspecto} & \textbf{Stored Procedures} & \textbf{Lógica en aplicación}\\
\midrule
Rendimiento & Excelente (compilados, cerca del dato) & Bueno con cache\\
Versionado con Git & Difícil (suele estar fuera del repo) & Natural\\
Testing & Difícil (requiere BD viva) & Fácil con mocks\\
Portabilidad SGBD & Acoplado a un motor & Multi-SGBD via ORM\\
Lógica compleja & Mantenimiento difícil & Más limpia\\
Operaciones masivas & Imbatible (sin round-trips) & Ineficiente\\
\bottomrule
\end{tabularx}
\end{cajacompara}

\textbf{Cuándo usar stored procedures en 2026:}
\begin{itemize}
\item Operaciones masivas que tocan miles de filas (ETL, batch).
\item Lógica que requiere consistencia transaccional fina con
múltiples tablas.
\item Reportes complejos con muchos JOINs.
\item Sistemas heredados donde la lógica ya está en SP.
\end{itemize}

\section{Práctica integrada: CRUD ORM vs SQL puro}

La práctica integrada de la semana confronta directamente
las dos estrategias dominantes: JDBC puro vs Spring Data JPA. El
estudiante implementa el mismo CRUD con ambas y produce una tabla
comparativa.


\begin{cajaejercicio}{Ejercicio 11.1 --- Comparación experimental SQL puro vs ORM (Sumativa GA-PE Unidad III)}

\textbf{Objetivo:} implementar el mismo CRUD en dos versiones (SQL
puro y ORM) y comparar líneas de código, rendimiento y mantenibilidad.

\subsubsection*{IDE: IntelliJ IDEA Ultimate}
Plugins necesarios (todos incluidos en Ultimate):
\begin{itemize}
\item \emph{Database Tools and SQL} (cliente BD nativo).
\item \emph{Spring Boot}.
\item \emph{Lombok}.
\end{itemize}

\subsubsection*{Paso 1 --- Preparar PostgreSQL en Docker}

Más simple que instalar PostgreSQL nativo (Listado~\ref{lst:11.15}):
\begin{lstlisting}[style=shell,numbers=none,caption={Ejemplo de Shell},label=lst:11.15]
docker run -d --name pg-appweb \ # arranca un contenedor
  -e POSTGRES_USER=appweb \
  -e POSTGRES_PASSWORD=appweb \
  -e POSTGRES_DB=appweb \
  -p 5432:5432 postgres:16
\end{lstlisting}

Verificar (Listado~\ref{lst:11.16}):
\begin{lstlisting}[style=shell,numbers=none,caption={Ejemplo de Shell},label=lst:11.16]
# ejecuta comando en contenedor activo
# ejecuta comando en contenedor activo
docker exec -it pg-appweb psql -U appweb -d appweb -c "SELECT version();"
\end{lstlisting}

\subsubsection*{Paso 2 --- Conectar IntelliJ a la BD}

Una vez levantada PostgreSQL, el siguiente paso es configurar
IntelliJ IDEA Ultimate para conectarse a ella desde la propia IDE.
Esto permite ejecutar SQL ad-hoc, ver datos y probar consultas sin
salir del entorno.


\begin{enumerate}
\item Abrir el panel \emph{Database} (lateral derecho).
\item Pulsar \texttt{+} $\rightarrow$ \texttt{Data Source}
$\rightarrow$ \texttt{PostgreSQL}.
\item Host: \texttt{localhost}, Port: \texttt{5432},
Database: \texttt{appweb}, User: \texttt{appweb}, Password:
\texttt{appweb}.
\item Pulsar \texttt{Test Connection}. Si pide descargar el driver,
aceptar.
\item Pulsar \texttt{OK}.
\end{enumerate}

\subsubsection*{Paso 3 --- Crear tabla y datos de prueba}

Click derecho sobre la conexión $\rightarrow$ \texttt{New}
$\rightarrow$ \texttt{Query Console}. Ejecutar (Listado~\ref{lst:11.17}):

\begin{lstlisting}[style=sql,numbers=none,caption={Ejemplo de SQL},label=lst:11.17]
-- crea tabla nueva en la BD
CREATE TABLE productos (                                    -- crea tabla nueva en la BD
    id        BIGSERIAL PRIMARY KEY,
    -- columna obligatoria (no admite NULL)
    nombre    VARCHAR(120) NOT NULL,                        -- columna obligatoria (no admite NULL)
    -- columna obligatoria (no admite NULL)
    precio    NUMERIC(10,2) NOT NULL CHECK (precio >= 0),   -- columna obligatoria (no admite NULL)
    -- columna obligatoria (no admite NULL)
    -- columna obligatoria (no admite NULL)
    stock     INTEGER NOT NULL DEFAULT 0 CHECK (stock >= 0),
    -- columna obligatoria (no admite NULL)
    creado    TIMESTAMPTZ NOT NULL DEFAULT NOW()            -- columna obligatoria (no admite NULL)
);

-- Cargar 1000 productos para benchmark
-- inserta una fila nueva
INSERT INTO productos (nombre, precio, stock)               -- inserta una fila nueva
SELECT
    'Producto ' || i,
    ROUND((random() * 100 + 1)::numeric, 2),
    (random() * 100)::int
-- tabla de origen
FROM generate_series(1, 1000) i;                            -- tabla de origen
\end{lstlisting}

Verificar: \texttt{SELECT COUNT(*) FROM productos;} debe devolver 1000.

\subsubsection*{Paso 4 --- Crear proyecto Spring Boot}

En IntelliJ: \texttt{File $\rightarrow$ New $\rightarrow$ Project
$\rightarrow$ Spring Boot}. Configurar:
\begin{itemize}
\item Language: Java; Build: Maven; JDK: 21.
\item Group: \texttt{ec.edu.uteq}; Artifact: \texttt{benchmark-orm}.
\item Spring Boot 4.x.
\item Dependencies: \emph{Spring Web}, \emph{Spring Data JPA},
\emph{PostgreSQL Driver}, \emph{Lombok}, \emph{Spring Boot DevTools}.
\end{itemize}

\subsubsection*{Paso 5 --- Configurar la conexión}

Editar \texttt{src/main/resources/application.yml} (Listado~\ref{lst:11.18}):

\begin{lstlisting}[style=yaml,numbers=none,caption={Ejemplo de YAML},label=lst:11.18]
spring: # configuracion de Spring
  datasource: # configuracion de origen de datos (BD)
    url:      jdbc:postgresql://localhost:5432/appweb
    username: appweb # usuario de la BD
    password: appweb # contrasena (mejor leer de env var)
  jpa: # configuracion JPA/Hibernate
    hibernate.ddl-auto: none # tabla ya existe
    show-sql: true # ver SQL generado
    properties.hibernate.format_sql: true
\end{lstlisting}

\subsubsection*{Paso 6 --- Versión A: JDBC puro}

La versión A implementa el CRUD usando solo JDBC puro: el
estudiante escribe el SQL a mano, gestiona la conexión, prepara las
sentencias y procesa los resultados. Es laborioso pero ilustrativo. V\'ease el Listado~\ref{lst:11.19}.


\begin{lstlisting}[style=java,caption={JdbcController.java},label=lst:11.19]
package ec.edu.uteq.benchmark.jdbc; // paquete del archivo

import org.springframework.jdbc.core.JdbcTemplate; // importacion de clase externa
import org.springframework.web.bind.annotation.*; // importacion de clase externa
import java.sql.Timestamp; // importacion de clase externa
import java.util.*; // importacion de clase externa

@RestController // expone endpoints REST con JSON
@RequestMapping("/api/jdbc/productos") // ruta base del controlador
public class JdbcController { // clase publica

    private final JdbcTemplate jdbc; // inmutable

    // metodo publico
    // metodo publico
    public JdbcController(JdbcTemplate jdbc) { this.jdbc = jdbc; }

    @GetMapping // endpoint HTTP GET
    public List<Map<String, Object>> listar() { // metodo publico
        return jdbc.queryForList( // retorna
            "SELECT id, nombre, precio, stock, creado FROM productos "
          + "ORDER BY id");
    }

    @GetMapping("/{id}") // endpoint HTTP GET
    // metodo publico
    // metodo publico
    public Map<String, Object> obtener(@PathVariable long id) {
        return jdbc.queryForMap( // retorna
            "SELECT id, nombre, precio, stock, creado FROM productos "
          + "WHERE id = ?", id);
    }
}
\end{lstlisting}

\subsubsection*{Paso 7 --- Versión B: Spring Data JPA}

La versión B implementa el mismo CRUD usando Spring Data
JPA: el estudiante define la entidad, declara la interfaz del
repositorio, y JPA genera el SQL automáticamente. El contraste con
la versión A es revelador. V\'ease el Listado~\ref{lst:11.20}.


\begin{lstlisting}[style=java,caption={Producto.java},label=lst:11.20]
package ec.edu.uteq.benchmark.jpa; // paquete del archivo

import jakarta.persistence.*; // importacion de clase externa
import lombok.*; // importacion de clase externa
import java.math.BigDecimal; // importacion de clase externa
import java.time.OffsetDateTime; // importacion de clase externa

@Entity @Table(name = "productos") // entidad JPA (mapea a tabla)
@Getter @Setter @NoArgsConstructor @AllArgsConstructor
public class Producto { // clase publica
    @Id @GeneratedValue(strategy = GenerationType.IDENTITY) // campo clave primaria
    private Long id; // campo privado
    private String nombre; // campo privado
    private BigDecimal precio; // campo privado
    private Integer stock; // campo privado
    private OffsetDateTime creado; // campo privado
}
\end{lstlisting}

\begin{lstlisting}[style=java,caption={ProductoRepository.java + JpaController.java},label=lst:11.21]
package ec.edu.uteq.benchmark.jpa; // paquete del archivo

// importacion de clase externa
// importacion de clase externa
import org.springframework.data.jpa.repository.JpaRepository;
public interface ProductoRepository // declaracion de interfaz
    extends JpaRepository<Producto, Long> {}

// ---

import org.springframework.web.bind.annotation.*; // importacion de clase externa
import java.util.*; // importacion de clase externa

@RestController // expone endpoints REST con JSON
@RequestMapping("/api/jpa/productos") // ruta base del controlador
public class JpaController { // clase publica
    private final ProductoRepository repo; // inmutable
    // metodo publico
    // metodo publico
    public JpaController(ProductoRepository r) { this.repo = r; }

    // endpoint HTTP GET
    // endpoint HTTP GET
    @GetMapping public List<Producto> listar() { return repo.findAll(); }

    @GetMapping("/{id}") // endpoint HTTP GET
    public Producto obtener(@PathVariable Long id) { // metodo publico
        return repo.findById(id).orElseThrow(); // retorna
    }
}
\end{lstlisting}

\subsubsection*{Paso 8 --- Benchmark con Apache Bench}

Ejecutar la app: \texttt{Shift+F10} sobre la clase principal. Servidor
en \url{http://localhost:8080}.

Desde otra terminal, instalar Apache Bench (\texttt{ab}). En Ubuntu:
\texttt{sudo apt install apache2-utils}; en macOS: viene preinstalado.

Ejecutar (Listado~\ref{lst:11.22}):
\begin{lstlisting}[style=shell,numbers=none,caption={Ejemplo de Shell},label=lst:11.22]
# 1000 peticiones, 10 concurrentes, version JDBC
ab -n 1000 -c 10 http://localhost:8080/api/jdbc/productos > jdbc.txt

# Misma prueba para JPA
ab -n 1000 -c 10 http://localhost:8080/api/jpa/productos > jpa.txt

cat jdbc.txt | grep "Requests per second"
cat jpa.txt  | grep "Requests per second"
\end{lstlisting}

\subsubsection*{Paso 9 --- Tabla comparativa para entregar}

Tras implementar las dos versiones, el estudiante debe
producir una tabla comparativa que cuantifique las diferencias.


\begin{tabularx}{\textwidth}{Xll}
\toprule
\textbf{Métrica} & \textbf{Versión JDBC} & \textbf{Versión JPA}\\
\midrule
Líneas de código del controller & & \\
Líneas de código del repository & & \\
Requests/segundo (1000 req, 10 conc) & & \\
Tiempo medio por request (ms) & & \\
SQL generado (consultas en log) & 1 fija & varía con caché L1/L2 \\
Vulnerabilidad SQL Injection & Eliminada (\texttt{?}) & Eliminada (JPA) \\
Mantenibilidad (1--5) & & \\
Productividad para CRUD nuevo (1--5) & & \\
\bottomrule
\end{tabularx}

\subsubsection*{Reflexión final escrita (entregable)}

Responder en máximo 300 palabras:
\begin{enumerate}
\item ¿Cuál fue más rápido en el benchmark?
\item ¿Qué versión escribiría más rápido un desarrollador nuevo?
\item ¿En qué casos el ORM es perjudicial (ineficiente)?
\item ¿Por qué muchas empresas mantienen \emph{ambos} estilos en el
mismo proyecto?
\end{enumerate}

\textbf{Esta es la \emph{Sumativa GA-PE Unidad III} de la semana 11.}
\end{cajaejercicio}

\section{Buenas prácticas profesionales en acceso a datos}

La caja siguiente compila las reglas profesionales del
acceso a datos: principios que separan el código mantenible del
código que se vuelve ingobernable a los seis meses.


\begin{cajabuenas}{Quince reglas profesionales 2026}
\begin{enumerate}
\item NUNCA concatenar datos del usuario en SQL: usar siempre
parámetros.
\item Configurar el driver con \texttt{ATTR\_EMULATE\_PREPARES = false}
en PHP.
\item DTOs separados de entidades: nunca exponer entidades JPA en
controllers.
\item Consultas paginadas siempre (\texttt{LIMIT} + \texttt{OFFSET})
para listados.
\item Resolver el problema N+1 con \emph{eager loading} explícito.
\item Migraciones con Flyway o Liquibase: jamás
\texttt{ddl-auto=update} en producción.
\item Índices en columnas de búsqueda frecuente y en foreign keys.
\item Conexiones desde un \emph{pool} (HikariCP en Spring, pgBouncer
externo).
\item Transacciones explícitas para operaciones multi-tabla.
\item Soft delete (\texttt{activo=false}) sobre delete físico cuando
hay relaciones.
\item Auditar cambios sensibles (\texttt{creado\_en},
\texttt{actualizado\_en}, \texttt{actualizado\_por}).
\item Cuenta BD con permisos mínimos: la app no necesita
\texttt{DROP TABLE}.
\item Encriptar la conexión con SSL/TLS en producción.
\item Tests con BD en memoria (H2) o testcontainers para integración
real.
\item Documentar el modelo entidad-relación en el repositorio
(ER en draw.io).
\end{enumerate}
\end{cajabuenas}

\section{Ejercicio resuelto adicional con resultado visual}

Como cierre de la semana 11, el siguiente ejercicio adicional
está completamente resuelto y se complementa con un propuesto.


\begin{cajaejemplo}{Ejercicio resuelto 11.1 --- CRUD de libros con Spring Data JPA y H2}
\textbf{Enunciado.} Implementar un CRUD completo de la entidad
\texttt{Libro} (id, titulo, autor, anio) con Spring Data JPA y H2
embebido, exponer los endpoints REST y verificar visualmente las
operaciones desde Postman.

\textbf{Entidad (\texttt{Libro.java}):}
\begin{lstlisting}[language=Java,caption={Entidad JPA con validaciones},label=lst:11.23]
@Entity // entidad JPA (mapea a tabla)
@Table(name = "libros") // nombre de la tabla en BD
public class Libro { // clase publica
  @Id @GeneratedValue(strategy = GenerationType.IDENTITY) // campo clave primaria
  private Long id; // campo privado
  @NotBlank @Size(max = 200) // validacion: no nulo y no vacio
  private String titulo; // campo privado
  @NotBlank @Size(max = 100) // validacion: no nulo y no vacio
  private String autor; // campo privado
  @Min(1900) @Max(2100) // validacion: valor minimo
  private int anio; // campo privado
  // getters/setters omitidos por brevedad
}
\end{lstlisting}

\textbf{Repositorio (\texttt{LibroRepository.java}):}
\begin{lstlisting}[language=Java,caption={Repositorio Spring Data},label=lst:11.24]
// declaracion de interfaz
// declaracion de interfaz
public interface LibroRepository extends JpaRepository<Libro, Long> {
  List<Libro> findByAutor(String autor);
  @Query("SELECT l FROM Libro l WHERE l.anio >= :desde") // JPQL
  List<Libro> publicadosDesde(@Param("desde") int desde);
}
\end{lstlisting}

\textbf{Controlador REST (\texttt{LibroController.java}):}
\begin{lstlisting}[language=Java,caption={Endpoints REST},label=lst:11.25]
@RestController @RequestMapping("/api/v1/libros") // expone endpoints REST con JSON
public class LibroController { // clase publica
  private final LibroRepository repo; // inmutable
  // metodo publico
  // metodo publico
  public LibroController(LibroRepository r) { this.repo = r; }

  // endpoint HTTP GET
  // endpoint HTTP GET
  @GetMapping public List<Libro> listar() { return repo.findAll(); }
  // endpoint HTTP GET
  // endpoint HTTP GET
  @GetMapping("/{id}") public ResponseEntity<Libro> obtener(@PathVariable Long id) {
    return repo.findById(id).map(ResponseEntity::ok) // retorna
      .orElse(ResponseEntity.notFound().build());
  }
  // endpoint HTTP POST (crear)
  // endpoint HTTP POST (crear)
  @PostMapping public ResponseEntity<Libro> crear(@Valid @RequestBody Libro l) {
    Libro guardado = repo.save(l);
    return ResponseEntity.status(201).body(guardado); // retorna
  }
  // endpoint HTTP PUT (actualizar completo)
  // endpoint HTTP PUT (actualizar completo)
  @PutMapping("/{id}") public ResponseEntity<Libro> actualizar(
      @PathVariable Long id, @Valid @RequestBody Libro l) { // extrae valor de la URL
    return repo.findById(id).map(existente -> { // retorna
      existente.setTitulo(l.getTitulo());
      existente.setAutor(l.getAutor());
      existente.setAnio(l.getAnio());
      return ResponseEntity.ok(repo.save(existente)); // retorna
    }).orElse(ResponseEntity.notFound().build());
  }
  // endpoint HTTP DELETE (eliminar)
  // endpoint HTTP DELETE (eliminar)
  @DeleteMapping("/{id}")
  public ResponseEntity<Void> eliminar(@PathVariable Long id) {
    return repo.findById(id).map(l -> { // retorna
      repo.delete(l);
      return ResponseEntity.noContent().<Void>build(); // retorna
    }).orElse(ResponseEntity.notFound().build());
  }
}
\end{lstlisting}

\textbf{Resultado visual} de \texttt{GET /api/v1/libros} en Postman
tras insertar tres libros:

\begin{center}
\fbox{\begin{minipage}{0.9\textwidth}
\vspace{0.2cm}\footnotesize
\hspace{0.3cm}\textbf{GET} \texttt{http://localhost:8080/api/v1/libros} $\quad$ \textbf{200 OK}\\[0.3cm]
\hspace{0.3cm}\texttt{[}\\
\hspace{0.6cm}\texttt{\{ \char34{}id\char34{}: 1, \char34{}titulo\char34{}: \char34{}Clean Code\char34{}, \char34{}autor\char34{}: \char34{}Robert C. Martin\char34{}, \char34{}anio\char34{}: 2008 \},}\\
\hspace{0.6cm}\texttt{\{ \char34{}id\char34{}: 2, \char34{}titulo\char34{}: \char34{}Designing Data-Intensive Applications\char34{},}\\
\hspace{1.0cm}\texttt{\char34{}autor\char34{}: \char34{}Martin Kleppmann\char34{}, \char34{}anio\char34{}: 2017 \},}\\
\hspace{0.6cm}\texttt{\{ \char34{}id\char34{}: 3, \char34{}titulo\char34{}: \char34{}Software Architecture in Practice\char34{},}\\
\hspace{1.0cm}\texttt{\char34{}autor\char34{}: \char34{}Len Bass\char34{}, \char34{}anio\char34{}: 2021 \}}\\
\hspace{0.3cm}\texttt{]}\\
\vspace{0.1cm}
\end{minipage}}
\end{center}

\textbf{Vista de la tabla H2} (consola H2 en \url{http://localhost:8080/h2-console}):

\begin{center}
\fbox{\begin{minipage}{0.9\textwidth}
\vspace{0.15cm}\footnotesize\centering
\begin{tabular}{|c|l|l|c|}
\hline
\textbf{ID} & \textbf{TITULO} & \textbf{AUTOR} & \textbf{ANIO} \\
\hline
1 & Clean Code & Robert C. Martin & 2008 \\
2 & Designing Data-Intensive Applications & Martin Kleppmann & 2017 \\
3 & Software Architecture in Practice & Len Bass & 2021 \\
\hline
\end{tabular}\\[0.1cm]
\end{minipage}}
\end{center}

\textbf{Discusión.} JPA generó automáticamente el SQL apropiado
(\texttt{INSERT}, \texttt{SELECT}, \texttt{UPDATE}, \texttt{DELETE})
y vinculó los parámetros por nombre, eliminando por completo la
posibilidad de SQL Injection. Las validaciones \texttt{@NotBlank},
\texttt{@Size} y \texttt{@Min}/\texttt{@Max} devuelven HTTP 400 con
detalle si el cliente envía datos inválidos.
\end{cajaejemplo}

\begin{cajaejercicio}{Ejercicio propuesto 11.2 --- Catálogo de cursos con relaciones y paginación}
\textbf{Enunciado.} Implementar un catálogo de cursos universitarios
con relación \emph{«un curso tiene muchos estudiantes»} y
\emph{«un estudiante puede estar en muchos cursos»} (\textbf{relación
muchos-a-muchos}). Resolver el problema N+1 mediante \texttt{JOIN
FETCH} y exponer endpoints REST con paginación, búsqueda y
ordenamiento.

\textbf{Entidades:}
\begin{itemize}
\item \texttt{Curso(id, codigo, nombre, creditos)}.
\item \texttt{Estudiante(id, cedula, nombre, email)}.
\item Tabla intermedia \texttt{inscripciones(curso\_id, estudiante\_id, fecha)}.
\end{itemize}

\textbf{Requisitos técnicos:}
\begin{enumerate}
\item Usar Spring Boot 4 + PostgreSQL 16 (no H2).
\item Mapear la relación muchos-a-muchos con \texttt{@ManyToMany}.
\item Endpoint \texttt{GET /api/v1/cursos?page=0\&size=10\&sort=nombre,asc}
debe usar \texttt{Pageable}.
\item Endpoint \texttt{GET /api/v1/cursos/\{id\}/estudiantes} debe
usar \texttt{JOIN FETCH} para evitar N+1.
\item Mostrar evidencia con Hibernate logs (\texttt{spring.jpa.show-sql=true})
de que se ejecuta UNA consulta en lugar de N+1.
\item Documentación OpenAPI con springdoc, accesible en
\url{/swagger-ui.html}.
\item Migraciones de BD con Flyway o Liquibase.
\item Cobertura JaCoCo $\geq 70\%$ con tests de integración usando
testcontainers.
\end{enumerate}

\textbf{Criterios de evaluación:}
\begin{itemize}
\item No hay SQL concatenado en ninguna parte del código.
\item El log muestra una sola consulta para listar cursos con sus
estudiantes (no N+1).
\item Validaciones \texttt{@Valid} correctas en todos los endpoints.
\item Documentación Swagger funcional.
\end{itemize}

\textbf{Lecturas recomendadas:} \cite{ottinger2017hibernate} para
optimización de queries; \cite{ottinger2017hibernate} para análisis del
problema N+1 en producción.
\end{cajaejercicio}

% =====================================================================
%       SEMANA 12: PATRONES ARQUITECTÓNICOS PARA APLICACIONES WEB
% =====================================================================
\chapter{Patrones de arquitectura para aplicaciones web}
\label{cap:semana12}

\epigraph{«La arquitectura es el conjunto de decisiones que son
difíciles de cambiar. Un buen arquitecto, por tanto, posterga las
decisiones innecesarias y razona profundamente las inevitables.»}
{--- \textit{Adaptado de Robert C.~Martin}}

\section*{Resultado de aprendizaje de la semana}
Al concluir la semana 12, el estudiante diseña aplicaciones web bajo
patrones arquitectónicos escalables (multicapa, hexagonal, SOA,
microservicios, EDA, P2P), documenta sus decisiones con el modelo C4
\cite{brown2023c4} y los Architectural Decision Records (ADR), y
evalúa los \emph{trade-offs} entre atributos de calidad ISO/IEC 25010
\cite{iso25010}.

\section{Principios fundacionales de arquitectura de software}

Antes de hablar de patrones concretos (capas, hexagonal,
microservicios), conviene fijar los principios fundacionales que los
gobiernan todos: separación de intereses, alta cohesión, bajo
acoplamiento, dependencias unidireccionales.


\begin{cajadefinicion}{¿Qué es la arquitectura de software?}
Según Bass, Clements y Kazman \cite{bass2021software}, la arquitectura
de un sistema computacional es \emph{el conjunto de estructuras
necesarias para razonar sobre el sistema, comprende los elementos del
software, las relaciones entre ellos y las propiedades de ambos}.

Tres ideas clave en esta definición:
\begin{itemize}
\item \textbf{Estructuras múltiples}: no hay UNA arquitectura, sino
varias vistas (estática, dinámica, despliegue).
\item \textbf{Razonamiento}: la arquitectura existe para que un
equipo pueda pensar el sistema sin perderse en detalles.
\item \textbf{Decisiones difíciles de cambiar}: lo arquitectónico es
lo que costaría reescribir si se cambiara.
\end{itemize}
\end{cajadefinicion}

\subsection{Los principios SOLID aplicados a la arquitectura}

Los principios SOLID, formulados originalmente por Robert C.
Martin para diseño orientado a objetos, se generalizan a la
arquitectura. La caja siguiente los traduce al contexto
arquitectónico.


\begin{table}[htbp]
\centering
\caption{SOLID, sus autores y significado en arquitectura.}
\footnotesize
\begin{tabularx}{\textwidth}{lXX}
\toprule
\textbf{Principio} & \textbf{Significado} & \textbf{Aplicación arquitectónica}\\
\midrule
S --- Single Responsibility & Una clase, una razón para cambiar & Cada módulo/servicio tiene un foco\\
O --- Open/Closed & Abierto a extensión, cerrado a modificación & Plugins, hooks, eventos\\
L --- Liskov Substitution & Subtipos sustituibles por sus tipos base & Polimorfismo correcto\\
I --- Interface Segregation & Interfaces pequeñas y enfocadas & APIs minimal por módulo\\
D --- Dependency Inversion & Depender de abstracciones, no concreciones & Inyección de dependencias\\
\bottomrule
\end{tabularx}
\end{table}

\subsection{Atributos de calidad ISO/IEC 25010}

El estándar ISO/IEC 25010 define las características de
calidad de un sistema software. Conocerlas permite razonar sobre
arquitectura de forma estructurada. La tabla siguiente las cataloga.


\begin{figure}[htbp]
\centering
\begin{tikzpicture}[
  every node/.style={font=\scriptsize,align=center},
  centro/.style={circle,draw=uteqverde,thick,fill=uteqverde!20,
                 minimum size=2.2cm,font=\small\bfseries},
  atrib/.style={rectangle,draw=uteqverde!80,fill=uteqverde!8,thick,
                rounded corners=4pt,minimum width=2cm,minimum height=0.7cm}]

  \node[centro] (q) {Calidad\\ISO 25010};

  \node[atrib,above=2cm of q] (func) {Idoneidad\\funcional};
  \node[atrib,right=2cm of func] (perf) {Eficiencia de\\desempeño};
  \node[atrib,right=of q] (comp) {Compatibilidad};
  \node[atrib,below right=2cm and 1cm of q] (usab) {Usabilidad};
  \node[atrib,below=2cm of q] (fiab) {Fiabilidad};
  \node[atrib,below left=2cm and 1cm of q] (segu) {Seguridad};
  \node[atrib,left=of q] (mant) {Mantenibilidad};
  \node[atrib,left=2cm of func] (port) {Portabilidad};

  \foreach \n in {func,perf,comp,usab,fiab,segu,mant,port}
    \draw[uteqverde!50,thick] (q) -- (\n);
\end{tikzpicture}
\caption{Los ocho atributos de calidad de la norma ISO/IEC 25010. Toda
decisión arquitectónica supone un \emph{trade-off} entre algunos de
estos atributos.}
\end{figure}

\section{Patrón 1: Arquitectura por capas (3-capas y 4-capas)}

La arquitectura por capas es el patrón más extendido en
aplicaciones empresariales: la aplicación se divide en niveles
horizontales con responsabilidades específicas y dependencias hacia
abajo.


\subsection{Estructura clásica}

La estructura clásica de tres capas (presentación, lógica,
datos) se ha extendido a cuatro al separar la lógica de aplicación
del dominio. La figura siguiente la ilustra.


\begin{figure}[htbp]
\centering
\begin{tikzpicture}[
  every node/.style={font=\footnotesize,align=center},
  capa/.style={rectangle,draw=uteqverde,thick,minimum width=10cm,
               minimum height=1.1cm,rounded corners=2pt}]
  \node[capa,fill=uteqverde!15]
    (pres) {\textbf{Presentación}\\
            Razor / Thymeleaf / Angular --- Controllers --- DTOs};
  \node[capa,fill=uteqverde!25,below=0.3cm of pres]
    (neg)  {\textbf{Negocio / Servicios}\\
            Casos de uso, reglas de dominio, validación};
  \node[capa,fill=uteqverde!35,below=0.3cm of neg]
    (dat)  {\textbf{Acceso a Datos / Repositorios}\\
            JPA / Doctrine / EF Core};
  \node[capa,fill=uteqverde!50,below=0.3cm of dat]
    (db)   {\textbf{Base de Datos}\\
            PostgreSQL / MySQL / SQL Server};

  \draw[->,>=stealth,thick] ([xshift=-3cm]pres.south)
       -- ([xshift=-3cm]neg.north);
  \draw[->,>=stealth,thick] ([xshift=-3cm]neg.south)
       -- ([xshift=-3cm]dat.north);
  \draw[->,>=stealth,thick] ([xshift=-3cm]dat.south)
       -- ([xshift=-3cm]db.north);
\end{tikzpicture}
\caption{Arquitectura clásica de 4 capas para aplicaciones web. La
dependencia siempre es descendente: la capa superior depende de la
inferior, nunca al revés.}
\end{figure}

\textbf{Regla cardinal:} una capa puede invocar a la inmediatamente
inferior; nunca al revés ni saltarse capas (no permitir que un
controller acceda directamente al SGBD).

\subsection{Implementación canónica en Spring Boot}

Spring Boot facilita la implementación de la arquitectura por
capas mediante sus anotaciones estereotipo
(\texttt{@Controller}, \texttt{@Service}, \texttt{@Repository}). El
ejemplo siguiente muestra una implementación canónica.


\begin{cajaejemplo}{Ejemplo 12.1 --- Estructura por capas en un proyecto Spring Boot}

\textbf{Enunciado.} Modelar arquitectónicamente con el modelo C4 (niveles 1 y 2) una tienda en línea sencilla. Producir el diagrama de Contexto y el de Contenedores.

\begin{lstlisting}[style=shell,numbers=none,caption={Ejemplo de Shell},label=lst:12.1]
src/main/java/ec/edu/uteq/appweb/
  controller/         <- Capa de Presentacion
    ProductoController.java
    ExceptionHandler.java
  service/            <- Capa de Negocio
    ProductoService.java
    impl/ProductoServiceImpl.java
  repository/         <- Capa de Acceso a Datos
    ProductoRepository.java
  model/              <- Entidades de dominio
    Producto.java
  dto/                <- Objetos de transferencia
    ProductoDto.java
    ProductoRequest.java
  exception/
    ResourceNotFoundException.java
    ValidationException.java
  config/
    SecurityConfig.java
    WebConfig.java
\end{lstlisting}
\end{cajaejemplo}

\subsection{Variante: Arquitectura Hexagonal (Ports and Adapters)}

La arquitectura hexagonal (\emph{Ports and Adapters}) de
Alistair Cockburn es una variante refinada de la arquitectura por
capas que coloca el dominio en el centro y vuelve invertibles las
dependencias.


\begin{figure}[htbp]
\centering
\begin{tikzpicture}[
  every node/.style={font=\scriptsize,align=center}]

  \node[regular polygon,regular polygon sides=6,
        draw=uteqverde,thick,fill=uteqverde!10,
        minimum width=4cm] (hex) {};
  \node[font=\small\bfseries,align=center] at (hex) {Dominio\\(Casos de uso)\\Puro Java};

  \node[rectangle,draw=uteqdorado,fill=uteqdorado!15,thick,
        rounded corners,minimum width=2cm,minimum height=0.8cm,
        above left=1cm and -0.5cm of hex] (rest) {REST\\Controller};
  \node[rectangle,draw=uteqdorado,fill=uteqdorado!15,thick,
        rounded corners,minimum width=2cm,minimum height=0.8cm,
        above right=1cm and -0.5cm of hex] (cli) {Cliente\\HTTP};
  \node[rectangle,draw=uteqdorado,fill=uteqdorado!15,thick,
        rounded corners,minimum width=2cm,minimum height=0.8cm,
        below left=1cm and -0.5cm of hex] (jpa) {JPA\\Repository};
  \node[rectangle,draw=uteqdorado,fill=uteqdorado!15,thick,
        rounded corners,minimum width=2cm,minimum height=0.8cm,
        below right=1cm and -0.5cm of hex] (cache) {Redis\\Cache};

  \draw[<->,thick] (rest)  -- (hex);
  \draw[<->,thick] (cli)   -- (hex);
  \draw[<->,thick] (jpa)   -- (hex);
  \draw[<->,thick] (cache) -- (hex);
\end{tikzpicture}
\caption{Arquitectura hexagonal: el dominio (centro) define
\emph{ports} (interfaces) y los \emph{adapters} (implementaciones)
externos lo conectan con tecnologías concretas (REST, JPA, Redis).
Permite cambiar el adaptador de persistencia sin tocar el dominio.}
\end{figure}

\section{Patrón 2: Arquitectura Orientada a Servicios (SOA)}

SOA emergió en los 2000 como enfoque para integrar sistemas
heterogéneos mediante servicios autónomos comunicados por SOAP/XML
sobre un \emph{Enterprise Service Bus} (ESB). Características:

\begin{itemize}
\item Servicios reutilizables expuestos vía contratos formales (WSDL).
\item Comunicación asíncrona vía ESB (TIBCO, MuleSoft, IBM IIB).
\item Gobernanza centralizada (registro de servicios, políticas
unificadas).
\item Datos a menudo compartidos entre servicios.
\end{itemize}

\textbf{Cuándo SOA sigue siendo relevante:} integraciones con
sistemas heredados (banca, gobierno), entornos donde el WSDL es
contractualmente exigido, dominios con altos requisitos de
interoperabilidad B2B.

\section{Patrón 3: Microservicios}

Microservicios es el patrón arquitectónico que más atención
ha recibido en la última década. Implica descomponer el sistema en
servicios pequeños, autónomos, desplegables y escalables
independientemente.


\begin{cajadefinicion}{Microservicios según Newman}
\emph{«Pequeños servicios autónomos que trabajan juntos, modelados
alrededor de un dominio de negocio.»} \cite{newman2021building}

\textbf{Características distintivas:}
\begin{itemize}
\item Servicios pequeños con responsabilidad acotada.
\item Despliegue independiente (cada servicio tiene su pipeline CI/CD).
\item Comunicación vía HTTP/JSON o eventos asíncronos.
\item Datos descentralizados (una BD por servicio).
\item Equipos autónomos (cada equipo es dueño de uno o varios servicios).
\item Heterogeneidad tecnológica permitida.
\end{itemize}
\end{cajadefinicion}

\subsection{Comparativa SOA vs microservicios}

SOA (\emph{Service-Oriented Architecture}) y microservicios
comparten la idea de modularidad por servicios, pero difieren en
escala, autonomía y tecnologías. La tabla siguiente los contrasta.


\begin{cajacompara}{SOA clásico vs microservicios modernos}
\centering\footnotesize
\begin{tabularx}{\textwidth}{lXX}
\toprule
\textbf{Criterio} & \textbf{SOA clásico} & \textbf{Microservicios}\\
\midrule
Tamaño del servicio & Mediano-grande & Pequeño y enfocado\\
Comunicación & SOAP/XML, ESB centralizado & REST/JSON, gRPC, eventos\\
Datos & Compartidos vía ESB & Una BD por servicio\\
Despliegue & Acoplado al ESB & Independiente por servicio\\
Equipos & Centralizados & Autónomos por servicio\\
Lenguajes & Generalmente Java/.NET & Heterogéneo\\
Coreografía/Orquestación & ESB orquesta & Coreografía con eventos\\
Madurez del enfoque & 2000s & 2015+\\
Complejidad operacional & Alta (ESB) & Alta (distribuida)\\
\bottomrule
\end{tabularx}
\end{cajacompara}

\subsection{Cuándo NO usar microservicios}

Microservicios no es la solución para todo problema. Hay
escenarios donde adoptarlos prematuramente genera más complejidad
que valor. La caja siguiente enumera los casos donde un monolito
modular es la mejor opción.


\begin{cajaadvertencia}{El error costoso de los microservicios prematuros}
Sam Newman y Martin Fowler son enfáticos: los microservicios son
una \textbf{arquitectura distribuida}, y los sistemas distribuidos
son \emph{inherentemente complejos}. Los problemas que un monolito
resuelve trivialmente (transacciones ACID, debugging,
deployment) se vuelven retos no triviales en microservicios.

\textbf{No use microservicios si:}
\begin{itemize}
\item Su equipo es pequeño (menos de 8 personas).
\item No tiene CI/CD maduro y observabilidad robusta.
\item No domina contenedores y orquestación (Kubernetes).
\item No ha experimentado los dolores reales de un monolito grande.
\end{itemize}

\textbf{La recomendación moderna} (Fowler 2024) es \emph{«monolith
first»}: empezar con un monolito modular bien diseñado, y migrar
servicios a microservicios solo cuando se justifique con datos.
\end{cajaadvertencia}

\subsection{Patrones esenciales de microservicios}

Los microservicios traen consigo problemas distribuidos
(latencia, fallos parciales, consistencia eventual) que se mitigan
con patrones específicos. La lista siguiente recoge los esenciales.


\begin{table}[htbp]
\centering
\caption{Patrones imprescindibles en arquitectura de microservicios.}
\footnotesize
\begin{tabularx}{\textwidth}{lX}
\toprule
\textbf{Patrón} & \textbf{Propósito}\\
\midrule
API Gateway & Punto único de entrada; auth, rate-limiting, routing.\\
Service Discovery & Registro y descubrimiento dinámico (Eureka, Consul).\\
Circuit Breaker & Cortar llamadas a servicios fallidos (Resilience4j).\\
Saga & Transacciones distribuidas vía secuencia de eventos.\\
CQRS & Separar comandos (escritura) de queries (lectura).\\
Event Sourcing & El estado se deriva de un log de eventos inmutable.\\
Sidecar & Funciones transversales en contenedor adyacente.\\
Bulkhead & Aislamiento de recursos para evitar cascada de fallos.\\
\bottomrule
\end{tabularx}
\end{table}

\section{Patrón 4: Arquitectura dirigida por eventos (EDA)}

En EDA los componentes se comunican mediante \emph{eventos}
(mensajes) en lugar de invocaciones directas. Tres roles:

\begin{itemize}
\item \textbf{Producer}: publica eventos cuando algo relevante ocurre.
\item \textbf{Broker}: encamina eventos (RabbitMQ, Apache Kafka, AWS
SNS/SQS).
\item \textbf{Consumer}: se suscribe y procesa eventos
asíncronamente.
\end{itemize}

\begin{figure}[htbp]
\centering
\begin{tikzpicture}[
  every node/.style={font=\scriptsize,align=center},
  caja/.style={rectangle,draw=uteqverde,thick,minimum width=2cm,
               minimum height=0.8cm,fill=uteqverde!10,rounded corners=2pt},
  flecha/.style={->,>=stealth,thick}]

  \node[caja] (p1) {Servicio\\Pedidos};
  \node[caja,below=of p1] (p2) {Servicio\\Pagos};

  \node[regular polygon,regular polygon sides=8,
        draw=uteqdorado,thick,fill=uteqdorado!15,
        minimum width=2.5cm,right=2cm of p1,yshift=-0.7cm] (br) {};
  \node[font=\bfseries\scriptsize] at (br) {Event\\Broker\\(Kafka)};

  \node[caja,right=2cm of br,yshift=0.8cm] (c1) {Servicio\\Email};
  \node[caja,right=2cm of br] (c2) {Servicio\\Inventario};
  \node[caja,right=2cm of br,yshift=-0.8cm] (c3) {Servicio\\Analitica};

  \draw[flecha] (p1) -- node[above]{publica} (br);
  \draw[flecha] (p2) -- node[below]{publica} (br);
  \draw[flecha] (br) -- node[above]{} (c1);
  \draw[flecha] (br) -- node[above]{} (c2);
  \draw[flecha] (br) -- node[above]{} (c3);
\end{tikzpicture}
\caption{Flujo en una arquitectura dirigida por eventos: los servicios
publican eventos sin saber quién los consumirá; los consumers se
suscriben a los eventos que les interesan.}
\end{figure}

\textbf{Ventajas:} desacoplamiento temporal y espacial,
escalabilidad horizontal trivial, resiliencia ante caídas parciales.

\textbf{Costos:} complejidad operacional, debugging distribuido
difícil, consistencia eventual (no ACID), \emph{message ordering}
problemático.

\section{Patrón 5: Peer-to-peer (P2P)}

Cada nodo es simultáneamente cliente y servidor. Histórica (Napster,
Gnutella, BitTorrent), actualmente vigente en \emph{blockchain},
\emph{IPFS}, dApps y comunicaciones cifradas (Briar, Jami).

\textbf{Para aplicaciones web empresariales tradicionales, el modelo
P2P rara vez es la respuesta.} Lo mencionamos porque el sílabo lo
incluye y porque puede aparecer si el PFC tiene componentes
descentralizados.

\section{Documentación arquitectónica con el modelo C4 de Brown}

El modelo C4 \cite{brown2023c4} estructura la documentación
arquitectónica en cuatro niveles de abstracción crecientes en detalle.

\subsection{Los cuatro niveles del modelo C4}

El modelo C4 de Simon Brown propone documentar la
arquitectura en cuatro niveles de zoom: Contexto, Contenedores,
Componentes y Código. Cada nivel responde preguntas distintas. La
descripción siguiente los detalla.


\begin{figure}[htbp]
\centering
\begin{tikzpicture}[
  every node/.style={font=\scriptsize,align=center},
  nivel/.style={rectangle,draw=uteqverde,thick,fill=uteqverde!10,
                rounded corners=4pt,minimum width=4cm,minimum height=1.5cm}]

  \node[nivel] (n1) {\textbf{Nivel 1: Contexto}\\
                     Cómo el sistema encaja\\con usuarios y sistemas externos};
  \node[nivel,right=0.5cm of n1] (n2)
       {\textbf{Nivel 2: Contenedores}\\
        Apps web, APIs, BDs,\\colas de mensajes};
  \node[nivel,below=0.6cm of n1] (n3)
       {\textbf{Nivel 3: Componentes}\\
        Descomposición interna\\de cada contenedor};
  \node[nivel,right=0.5cm of n3] (n4)
       {\textbf{Nivel 4: Código}\\
        Diagrama de clases UML,\\opcional};

  \draw[->,>=stealth,thick,uteqverde!70!black] (n1) -- (n2);
  \draw[->,>=stealth,thick,uteqverde!70!black] (n2) -- (n3);
  \draw[->,>=stealth,thick,uteqverde!70!black] (n3) -- (n4);
\end{tikzpicture}
\caption{Los cuatro niveles del modelo C4. Cada nivel hace zoom sobre
un elemento del nivel superior. La mayoría de equipos solo necesita
los niveles 1, 2 y 3; el nivel 4 se genera automáticamente desde el
código.}
\end{figure}

\subsection{Ejemplo: diagrama de contexto del PFC}

El diagrama de contexto del PFC muestra el sistema como una
caja única y sus interacciones con usuarios y sistemas externos. Es
el primer artefacto que debe acompañar todo PFC.


\begin{figure}[htbp]
\centering
\begin{tikzpicture}[node distance=1.5cm and 2cm]
  \cfourPerson{est}{(0,0)}{Estudiante}
  \cfourPerson{adm}{(0,-2.5)}{Administrador}
  \cfourSystem{sis}{(5,-1.25)}{Sistema PFC\\(API + SPA Angular)}
  \cfourSystem{smtp}{(10,0)}{Servicio SMTP\\externo (SendGrid)}
  \cfourSystem{idp}{(10,-2.5)}{Identidad UTEQ\\(LDAP)}

  \draw[->,thick,>=stealth] (est) -- node[above,font=\tiny]{HTTPS} (sis);
  \draw[->,thick,>=stealth] (adm) -- node[below,font=\tiny]{HTTPS} (sis);
  \draw[->,thick,>=stealth] (sis) -- node[above,font=\tiny]{SMTP} (smtp);
  \draw[->,thick,>=stealth] (sis) -- node[above,font=\tiny]{LDAP} (idp);
\end{tikzpicture}
\caption{Nivel 1 (Contexto) del modelo C4 para el PFC: usuarios y
sistemas externos.}
\end{figure}

\subsection{Ejemplo: diagrama de contenedores del PFC}

El diagrama de contenedores zoom-in sobre la caja del nivel
1 y revela las piezas tecnológicas mayores: frontend SPA, backend
API, base de datos, cache, servicios externos. La figura siguiente
lo ilustra para el PFC.


\begin{figure}[htbp]
\centering
\begin{tikzpicture}[node distance=1.5cm and 1.8cm]
  \cfourPerson{usr}{(0,0)}{Usuario\\(Estudiante)}

  \cfourContainer{spa}{(4,0)}{Frontend\\Angular 19\\(SPA)}
  \cfourContainer{api}{(8,0)}{API REST\\Spring Boot 4\\(Tomcat 10.1)}
  \cfourDb{pg}{(12,0)}{PostgreSQL\\16}
  \cfourDb{rd}{(12,-2.5)}{Redis 7.4\\(cache)}

  \draw[->,thick,>=stealth] (usr.east) -- node[above,font=\tiny]{HTTPS} (spa.west);
  \draw[->,thick,>=stealth] (spa.east) -- node[above,font=\tiny]{JSON} (api.west);
  \draw[->,thick,>=stealth] (api.east) -- node[above,font=\tiny]{JDBC/SSL} (pg.west);
  \draw[->,thick,>=stealth] (api) -- node[right,font=\tiny]{TCP} (rd);
\end{tikzpicture}
\caption{Nivel 2 (Contenedores) del PFC: las unidades desplegables
y la tecnología de cada una.}
\end{figure}

\section{Architectural Decision Records (ADR)}

Un ADR documenta una decisión arquitectónica significativa.
Formato canónico (Michael Nygard, 2011) consta de cuatro secciones:
\textbf{Contexto}, \textbf{Decisión}, \textbf{Estado},
\textbf{Consecuencias}.

\begin{cajaadr}{ADR-001 --- Adoptar JWT en lugar de sesiones de servidor}
\textbf{Estado:} Aceptada el 2026-05-04.

\textbf{Contexto:} El frontend del PFC es una SPA Angular y el
backend es una API REST que requiere escalar horizontalmente con
Docker Compose en el laboratorio y, eventualmente, en cloud. Las
sesiones en memoria del servidor obligarían a \emph{sticky sessions}
en el balanceador o a un \emph{session store} compartido (Redis,
JDBC), lo que añade complejidad operacional.

\textbf{Decisión:} Adoptar autenticación con JSON Web Token (JWT)
firmado HS256 con secreto de 256 bits, transportado en cookie
\texttt{HttpOnly} \texttt{Secure} \texttt{SameSite=Strict}, con
\emph{refresh token} de 7 días y \emph{blacklist} en Redis para
revocación inmediata por JTI.

\textbf{Consecuencias:}
\begin{itemize}
\item[(+)] Backend stateless: cualquier instancia puede validar.
\item[(+)] Compatible con CORS sin sticky sessions.
\item[(+)] Incluye claims (rol, email) sin consulta extra a BD.
\item[(--)] Revocación requiere blacklist (mayor complejidad).
\item[(--)] El JWT es más grande que un session ID (\textasciitilde 350 bytes vs 32).
\item[(--)] Si se compromete el secreto, todos los tokens son
inválidos.
\end{itemize}

\textbf{Alternativas consideradas:}
\begin{itemize}
\item Sesiones server-side con Spring Session + Redis: descartada
por complejidad operacional desproporcionada.
\item OAuth 2.0 con servidor de autorización dedicado (Keycloak):
descartada por excesiva infraestructura para un PFC.
\end{itemize}
\end{cajaadr}

\section{Práctica integrada: documentar el PFC con C4 y ADRs}

La práctica de la semana es documentar el PFC del equipo con
diagramas C4 (niveles 1 y 2) y al menos dos ADRs justificando las
decisiones arquitectónicas mayores.


\begin{cajaejercicio}{Ejercicio 12.1 --- Entrega 2 del PFC --- Arquitectura Completa (Sumativa)}


\textbf{Enunciado.} Producir la Entrega 2 completa del PFC del equipo: arquitectura documentada con C4, API REST funcional con OpenAPI y frontend Angular consumiéndola. Sumativa principal del segundo corte.

\textbf{Esta es la \emph{Sumativa TA --- PFC Entrega 2: Arquitectura
Completa + API REST Documentada + Frontend Angular Funcional}.}

\subsubsection*{IDE y herramientas}

Antes de los pasos, conviene confirmar las herramientas que
el estudiante debe tener instaladas para la práctica.


\begin{itemize}
\item \textbf{IntelliJ IDEA Ultimate} (con plugins Spring Boot,
Database, Docker).
\item \textbf{Structurizr DSL} en línea (\url{https://www.structurizr.com/dsl})
o \textbf{draw.io} integrado en IntelliJ
(plugin \emph{Diagrams.net Integration}).
\item \textbf{Postman} para probar endpoints.
\end{itemize}

\subsubsection*{Paso 1 --- Generar diagrama C4 Nivel 1 (Contexto)}

En \url{https://www.structurizr.com/dsl}, crear un workspace y pegar (Listado~\ref{lst:12.2}):

\begin{lstlisting}[style=shell,numbers=none,caption={Ejemplo de Shell},label=lst:12.2]
workspace "PFC UTEQ" {
  model {
    estudiante = person "Estudiante"
    admin      = person "Administrador"
    sistema    = softwareSystem "Sistema PFC" {
      tags "App Principal"
    }
    smtp = softwareSystem "Servicio SMTP" {
      tags "Externo"
    }
    estudiante -> sistema "Usa via HTTPS"
    admin      -> sistema "Administra"
    sistema    -> smtp    "Envia notificaciones"
  }
  views {
    systemContext sistema "Contexto" {
      include *
      autoLayout
    }
  }
}
\end{lstlisting}

Pulsar \texttt{Render}: aparece el diagrama. Exportar como PNG.

\subsubsection*{Paso 2 --- Generar diagrama C4 Nivel 2 (Contenedores)}

Extender el modelo (Listado~\ref{lst:12.3}):

\begin{lstlisting}[style=shell,numbers=none,caption={Ejemplo de Shell},label=lst:12.3]
sistema = softwareSystem "Sistema PFC" {
  spa     = container "SPA Angular"     "Angular 19" "TypeScript"
  api     = container "API REST"        "Spring Boot 4" "Java"
  db      = container "Base de datos"   "PostgreSQL 16"  "BD" {
    tags "Database"
  }
  cache   = container "Cache"           "Redis 7.4"      "BD" {
    tags "Database"
  }
  spa -> api "JSON via HTTPS"
  api -> db  "JDBC/SSL"
  api -> cache "TCP"
}
\end{lstlisting}

\subsubsection*{Paso 3 --- Crear los 5 ADRs en el repositorio Git}

En el repositorio del PFC, crear directorio \texttt{docs/adr/}.
Crear estos archivos en formato Markdown:

\begin{itemize}
\item \texttt{ADR-001-JWT-vs-sesiones.md}
\item \texttt{ADR-002-PostgreSQL-vs-MySQL.md}
\item \texttt{ADR-003-Spring-Boot-vs-Quarkus.md}
\item \texttt{ADR-004-Redis-cache-aside.md}
\item \texttt{ADR-005-Docker-Compose-deployment.md}
\end{itemize}

Cada ADR sigue la plantilla de Nygard mostrada en este capítulo.

\subsubsection*{Paso 4 --- Tabla de atributos de calidad ISO 25010}

Documentar en el informe técnico esta tabla:

\begin{tabularx}{\textwidth}{llXX}
\toprule
\textbf{Atributo} & \textbf{Prioridad} & \textbf{Escenario} & \textbf{Estrategia}\\
\midrule
Seguridad & Crítica & SQL Injection en endpoints & JPA + Bean Validation\\
Eficiencia & Alta & Listado de productos < 200ms & Cache Redis con TTL 60s\\
Usabilidad & Alta & Login intuitivo & Angular con feedback visual\\
Mantenibilidad & Alta & Nuevo dev productivo en 1 día & C4 + ADRs + README\\
Fiabilidad & Media & Disponibilidad 99\% & Health checks + Docker restart\\
Portabilidad & Media & Mover de Windows a Linux & Docker Compose\\
\bottomrule
\end{tabularx}

\subsubsection*{Paso 5 --- API REST documentada con Swagger UI}

Agregar al \texttt{pom.xml} (Listado~\ref{lst:12.4}):

\begin{lstlisting}[style=xml,numbers=none,caption={Ejemplo de XML},label=lst:12.4]
<dependency>
  <groupId>org.springdoc</groupId>
  <artifactId>springdoc-openapi-starter-webmvc-ui</artifactId>
  <version>2.6.0</version>
</dependency>
\end{lstlisting}

Reiniciar la aplicación. Abrir
\url{http://localhost:8080/swagger-ui.html}: aparece toda la API
documentada y probable interactivamente.

\subsubsection*{Paso 6 --- Frontend Angular conectado}

Crear el frontend con Angular CLI (Listado~\ref{lst:12.5}):
\begin{lstlisting}[style=shell,numbers=none,caption={Ejemplo de Shell},label=lst:12.5]
ng new frontend --standalone --routing --style=scss
cd frontend # cambia directorio
ng generate service auth
ng generate component login --standalone
\end{lstlisting}

Implementar el flujo de login que llame al endpoint
\texttt{POST /api/auth/login} y guarde el JWT en cookie HttpOnly
(que el backend establece).

\subsubsection*{Paso 7 --- Empaquetar todo con Docker Compose}

Crear \texttt{docker-compose.yml} en la raíz (Listado~\ref{lst:12.6}):

\begin{lstlisting}[style=yaml,numbers=none,caption={Ejemplo de YAML},label=lst:12.6]
services:
  postgres:
    image: postgres:16-alpine
    environment:
      POSTGRES_USER: appweb
      POSTGRES_PASSWORD: appweb
      POSTGRES_DB:   appweb
    ports: ["5432:5432"]
    volumes: ["pgdata:/var/lib/postgresql/data"]
    healthcheck:
      test: ["CMD","pg_isready","-U","appweb"]
      interval: 10s

  redis: # configuracion del cliente Redis
    image: redis:7-alpine # configuracion del cliente Redis
    ports: ["6379:6379"]

  api:
    build: ./backend
    depends_on:
      postgres: { condition: service_healthy }
      redis:    { condition: service_started } # configuracion del cliente Redis
    environment:
      SPRING_DATASOURCE_URL: jdbc:postgresql://postgres:5432/appweb
      SPRING_REDIS_HOST: redis
    ports: ["8080:8080"]

  web:
    build: ./frontend
    depends_on: [api]
    ports: ["80:80"]

volumes:
  pgdata:
\end{lstlisting}

\subsubsection*{Paso 8 --- Verificar el despliegue completo}

Como verificación final, el estudiante debe levantar el
sistema completo y validar que funciona end-to-end antes de
considerar terminada la entrega. V\'ease el Listado~\ref{lst:12.7}.


\begin{lstlisting}[style=shell,numbers=none,caption={Ejemplo de Shell},label=lst:12.7]
docker compose up -d # levanta los servicios del compose
docker compose ps # todos healthy
curl http://localhost:8080/actuator/health
# Abrir http://localhost en el navegador: SPA cargada
\end{lstlisting}

\subsubsection*{Entregables}

La entrega de la práctica de la semana 12 debe incluir los
siguientes artefactos consolidados.


\begin{enumerate}
\item Repositorio Git con tag \texttt{v0.5.0}.
\item Documento PDF (15--25 páginas) con: portada UTEQ, diagrama C4
nivel 1 y 2 (PNG), tabla atributos calidad, los 5 ADRs, capturas de
Swagger UI funcional, capturas de la SPA Angular en login,
\texttt{docker-compose.yml} comentado.
\item Demostración en clase: \texttt{docker compose up} arranca todo
en menos de 2 minutos.
\end{enumerate}

\textbf{Esta es la \emph{Sumativa de la semana 12}: PFC Entrega 2.}
\end{cajaejercicio}

\section{Buenas prácticas profesionales en arquitectura}

La caja siguiente compila quince reglas profesionales de
arquitectura que el estudiante debe respetar al diseñar cualquier
sistema.


\begin{cajabuenas}{Doce reglas profesionales 2026}
\begin{enumerate}
\item \emph{Monolith first}: empezar con monolito modular bien
diseñado.
\item Documentar TODA decisión significativa con un ADR.
\item Diagramas C4 nivel 1 y 2 obligatorios; nivel 3 si el sistema
es grande.
\item Atributos de calidad explícitos al inicio del proyecto.
\item Versionar la documentación arquitectónica como código (Git).
\item Revisar la arquitectura cada trimestre con todo el equipo.
\item Evitar acoplamiento entre servicios vía base de datos
compartida.
\item Cada servicio dueño de su esquema (en microservicios).
\item Comunicación asíncrona cuando se pueda; síncrona solo cuando
sea inevitable.
\item Health checks y observabilidad desde el día 1.
\item Tests de arquitectura con ArchUnit (Java) o NetArchTest (.NET).
\item Pensar en \emph{evolutionary architecture}: cambiar es la
norma, no la excepción.
\end{enumerate}
\end{cajabuenas}

\section{Ejercicio resuelto adicional con resultado visual}

Como cierre, el siguiente ejercicio resuelto adicional
incluye un caso completo con diagramas y ADR; lo acompaña un
propuesto.


\begin{cajaejemplo}{Ejercicio resuelto 12.1 --- Diagrama C4 nivel 1 y 2 de una tienda en línea}
\textbf{Enunciado.} Modelar arquitectónicamente con el modelo C4 una
tienda en línea sencilla que tiene: SPA Angular para clientes,
backend Spring Boot REST, PostgreSQL, Redis para caché de catálogo y
servicio externo de pagos (Stripe). Producir los diagramas de nivel 1
(Sistema-Contexto) y nivel 2 (Contenedores).

\textbf{Nivel 1 --- Diagrama de Contexto del Sistema}
(representación textual estilo ASCII; en producción se renderiza
con PlantUML o Structurizr):

\begin{center}
\fbox{\begin{minipage}{0.9\textwidth}
\vspace{0.2cm}\centering\footnotesize
\begin{tabular}{c}
\fbox{\textbf{Cliente}}\\[0.1cm]
$\downarrow$\\[0.1cm]
\fbox{\textbf{Sistema TiendaUTEQ}}\\[0.1cm]
$\downarrow$ \\[0.1cm]
\fbox{\textbf{Stripe API}}\\
\end{tabular}\\[0.2cm]
\textit{El cliente interactúa con TiendaUTEQ; el sistema consume Stripe para pagos.}
\vspace{0.1cm}
\end{minipage}}
\end{center}

\textbf{Nivel 2 --- Diagrama de Contenedores:}

\begin{figure}[H]
\centering
\begin{tikzpicture}[
  every node/.style={font=\footnotesize,align=center},
  persona/.style={circle,draw=uteqverde,thick,fill=uteqverde!15,minimum size=1.4cm},
  contenedor/.style={rectangle,draw=uteqverde,thick,fill=uteqverde!10,
                     rounded corners=4pt,minimum width=2.8cm,minimum height=1.2cm},
  externo/.style={rectangle,draw=uteqdorado,thick,fill=uteqdorado!15,
                  dashed,minimum width=2.8cm,minimum height=1.2cm}]
  \node[persona] (u) {Cliente};
  \node[contenedor,right=1cm of u] (spa) {SPA Angular\\(Browser)};
  \node[contenedor,right=1cm of spa] (api) {API REST\\(Spring Boot)};
  \node[contenedor,below=0.6cm of api] (db) {BD\\(PostgreSQL)};
  \node[contenedor,left=0.6cm of db] (redis) {Cache\\(Redis)};
  \node[externo,right=1cm of api] (stripe) {Stripe};

  \draw[->,thick] (u) -- node[above,font=\tiny]{HTTPS} (spa);
  \draw[->,thick] (spa) -- node[above,font=\tiny]{JSON/HTTPS} (api);
  \draw[->,thick] (api) -- node[right,font=\tiny]{JDBC} (db);
  \draw[->,thick] (api) -- node[above,font=\tiny]{TCP} (redis);
  \draw[->,thick] (api) -- node[above,font=\tiny]{HTTPS} (stripe);
\end{tikzpicture}
\caption{Diagrama C4 nivel 2 (contenedores) de TiendaUTEQ.}
\end{figure}

\textbf{ADR-001 correspondiente:}

\begin{cajaadr}{ADR-001 --- Adopción de arquitectura por capas con caché Redis}
\textbf{Estado:} Aceptado el 5 de mayo de 2026.

\textbf{Contexto.} TiendaUTEQ requiere servir hasta 500
usuarios concurrentes con tiempos de respuesta $<$ 300 ms en el
listado de productos. La BD PostgreSQL muestra latencias de 50--80 ms
en consultas frecuentes.

\textbf{Decisión.} Adoptar arquitectura por capas (presentación,
aplicación, dominio, infraestructura) con Spring Boot y agregar caché
Redis al servicio de catálogo con TTL de 60 segundos.

\textbf{Consecuencias positivas:}
\begin{itemize}
\item Tiempos de respuesta $<$ 50 ms tras el primer hit.
\item Menor carga sobre PostgreSQL.
\end{itemize}

\textbf{Consecuencias negativas:}
\begin{itemize}
\item Complejidad operativa adicional (un contenedor más en Docker).
\item Posibilidad de inconsistencia transitoria.
\end{itemize}

\textbf{Alternativas consideradas:}
\begin{itemize}
\item \emph{Caffeine in-process}: descartado por no compartir caché
entre instancias.
\item \emph{Memcached}: descartado por inferior soporte de
estructuras de datos respecto a Redis.
\end{itemize}
\end{cajaadr}

\textbf{Discusión.} El estudiante debe notar que el diagrama no
muestra clases ni código: el nivel arquitectónico se ocupa de las
piezas grandes y las decisiones difíciles de revertir
\cite{bass2021software}. El ADR documenta la decisión para futuros
mantenedores que se pregunten \emph{«¿por qué Redis?»} en lugar de
adivinar.
\end{cajaejemplo}

\begin{cajaejercicio}{Ejercicio propuesto 12.2 --- Migrar un monolito a arquitectura hexagonal}
\textbf{Enunciado.} Tomar el PFC del equipo (o el monolito tienda en
línea anterior) e identificar los \emph{puertos} y \emph{adaptadores}
para refactorizarlo a arquitectura hexagonal (\emph{Ports and
Adapters} de Alistair Cockburn).

\textbf{Procedimiento:}
\begin{enumerate}
\item Identificar el \textbf{dominio puro}: entidades y servicios
de negocio sin dependencia de Spring, JPA, REST, etc.
\item Definir \textbf{puertos primarios} (interfaces que el dominio
expone para que lo invoquen): casos de uso.
\item Definir \textbf{puertos secundarios} (interfaces que el dominio
necesita): repositorios, servicios externos, notificación.
\item Crear \textbf{adaptadores primarios}: controladores REST,
listeners de eventos, CLI.
\item Crear \textbf{adaptadores secundarios}: implementaciones JPA
de repositorios, clientes HTTP de servicios externos, productores de
eventos.
\item Reorganizar el proyecto en Maven multi-módulo:
\texttt{dominio}, \texttt{aplicacion}, \texttt{adaptadores-rest},
\texttt{adaptadores-persistencia}.
\item Producir test del dominio que NO requieran arrancar Spring
(unit tests puros, $<$ 100 ms cada uno).
\end{enumerate}

\textbf{Criterios de evaluación:}
\begin{itemize}
\item El módulo \texttt{dominio} no tiene ninguna dependencia a
Spring, JPA ni infraestructura.
\item Existen tests del dominio que se ejecutan sin Spring
(\texttt{@SpringBootTest} no aparece).
\item Hay al menos dos adaptadores secundarios alternativos
(ej. \texttt{InMemoryUsuarioRepository} para tests +
\texttt{JpaUsuarioRepository} para producción).
\item Documentación con diagrama de la arquitectura y ADR-002 con la
decisión.
\end{itemize}

\textbf{Lecturas recomendadas:} \cite{ford2017building} para
\emph{evolutionary architecture}; \cite{newman2021building} sobre
descomposición incremental.
\end{cajaejercicio}

% =====================================================================
%       SEMANA 13: DISEÑO DE ARQUITECTURAS WEB ESCALABLES
% =====================================================================
\chapter{Diseño de arquitecturas web escalables}
\label{cap:semana13}

\epigraph{«Hay solo dos cosas difíciles en informática: invalidar la
caché y nombrar las cosas.»}{--- \textit{Phil Karlton}}

\section*{Resultado de aprendizaje de la semana}
Al concluir la semana 13, el estudiante diseña arquitecturas web
escalables aplicando estrategias de balanceo de carga, patrones de
caché distribuida con Redis, técnicas de evaluación arquitectónica
(ATAM) y comprende los \emph{trade-offs} entre escalabilidad,
consistencia y disponibilidad (teorema CAP).

\section{Escalabilidad en aplicaciones web: vertical vs horizontal}

Cuando una aplicación gana usuarios, llega el momento de
escalar. Hay dos estrategias fundamentales, vertical y horizontal,
con implicaciones técnicas y económicas muy diferentes.


\subsection{Definiciones formales}

La distinción entre escalado vertical y horizontal es la
base del razonamiento sobre arquitectura escalable. Definiciones
formales y casos típicos a continuación.


\begin{cajadefinicion}{Las dos dimensiones de la escalabilidad}
\textbf{Escalabilidad vertical} (\emph{scale up}): aumentar la
capacidad del mismo servidor con más CPU, más RAM, más disco. Se
detiene en algún momento (Ley de Moore agotada, costos exponenciales).

\textbf{Escalabilidad horizontal} (\emph{scale out}): añadir más
servidores idénticos detrás de un balanceador de carga. Crece
linealmente hasta el límite del componente menos escalable
(típicamente la BD).
\end{cajadefinicion}

\begin{table}[htbp]
\centering
\caption{Comparativa: escalado vertical vs horizontal.}
\footnotesize
\begin{tabularx}{\textwidth}{lXX}
\toprule
\textbf{Aspecto} & \textbf{Vertical} & \textbf{Horizontal}\\
\midrule
Costo inicial & Bajo (un servidor mayor) & Alto (infraestructura, balanceador)\\
Costo a gran escala & Exponencial & Lineal\\
Complejidad arquitectónica & Mínima & Alta (sesiones, consistencia)\\
Tolerancia a fallos & Baja (SPOF) & Alta (réplicas)\\
Aplicabilidad a BD & Limitada & Sharding requerido\\
Tiempo para añadir capacidad & Minutos (en cloud) & Segundos (autoscaling)\\
Lím. teórico & Capacidad máxima del hw & Prácticamente ilimitado\\
\bottomrule
\end{tabularx}
\end{table}

\subsection{El teorema CAP de Brewer}

El teorema CAP de Eric Brewer ---imposibilidad de tener
consistencia, disponibilidad y tolerancia a particiones
simultáneamente--- es uno de los resultados más citados en
arquitectura distribuida. Su interpretación práctica es más matizada
de lo que parece.


\begin{cajadefinicion}{Teorema CAP --- Eric Brewer 2000}
En un sistema distribuido, ante un \textbf{P}article \textbf{P}artition
(falla de red), un sistema solo puede garantizar
\emph{a lo sumo} dos de las tres propiedades siguientes:

\begin{description}[leftmargin=1.5cm]
\item[Consistency (C).] Toda lectura recibe el dato más reciente.
\item[Availability (A).] Toda petición recibe respuesta (no errores).
\item[Partition tolerance (P).] El sistema sigue funcionando ante
fallos de red entre nodos.
\end{description}

Como las particiones de red en sistemas distribuidos son inevitables,
en la práctica se elige entre \textbf{CP} (consistencia sobre
disponibilidad: PostgreSQL replicado) o \textbf{AP} (disponibilidad
sobre consistencia, con consistencia eventual: Cassandra, DynamoDB).
\end{cajadefinicion}

\section{Estrategias de balanceo de carga}

Un balanceador de carga distribuye las peticiones entrantes
entre múltiples instancias del backend. Es la pieza fundamental para
escalar horizontalmente.


\begin{figure}[htbp]
\centering
\begin{tikzpicture}[
  every node/.style={font=\scriptsize,align=center},
  caja/.style={rectangle,draw=uteqverde,thick,minimum width=1.8cm,
               minimum height=0.8cm,fill=uteqverde!10,rounded corners=2pt}]

  \node[caja] (cli) {Cliente};
  \node[caja,right=1.5cm of cli] (lb) {Load\\Balancer};

  \node[caja,right=1.5cm of lb,yshift=1.5cm] (a1) {App\\Server 1};
  \node[caja,right=1.5cm of lb] (a2) {App\\Server 2};
  \node[caja,right=1.5cm of lb,yshift=-1.5cm] (a3) {App\\Server 3};

  \node[cylinder,draw=uteqverde,fill=uteqverde!15,thick,
        shape border rotate=90,aspect=0.3,minimum width=1.5cm,
        minimum height=1.5cm,right=1.5cm of a2] (db) {BD};

  \draw[->,thick,>=stealth] (cli) -- (lb);
  \draw[->,thick,>=stealth] (lb) -- (a1);
  \draw[->,thick,>=stealth] (lb) -- (a2);
  \draw[->,thick,>=stealth] (lb) -- (a3);
  \draw[->,thick,>=stealth] (a1.east) -- (db.west);
  \draw[->,thick,>=stealth] (a2) -- (db);
  \draw[->,thick,>=stealth] (a3.east) -- (db.west);
\end{tikzpicture}
\caption{Topología clásica con balanceador de carga: el LB recibe las
peticiones del cliente y las distribuye entre N servidores de
aplicación que comparten la BD.}
\end{figure}

\subsection{Algoritmos de balanceo}

Distintos algoritmos de balanceo tienen comportamientos
diferentes ante distintas cargas. La tabla siguiente los cataloga.


\begin{table}[htbp]
\centering
\caption{Algoritmos de balanceo de carga.}
\footnotesize
\begin{tabularx}{\textwidth}{lXX}
\toprule
\textbf{Algoritmo} & \textbf{Descripción} & \textbf{Mejor uso}\\
\midrule
Round-robin & Distribución uniforme rotativa & Servidores homogéneos\\
Least connections & Al servidor con menos conexiones activas & Carga heterogénea\\
Least response time & Al más rápido & APIs con latencia variable\\
IP hash & Hash IP $\to$ servidor (sticky) & Sesiones en memoria\\
Weighted round-robin & RR con pesos por servidor & Servidores asimétricos\\
Random & Aleatorio uniforme & Pruebas, casos especiales\\
\bottomrule
\end{tabularx}
\end{table}

\subsection{Implementaciones populares}

\textbf{Nginx} es el balanceador más usado para aplicaciones web:

\begin{cajaejemplo}{Ejemplo 13.1 --- Nginx como load balancer}

\textbf{Enunciado.} Implementar el patrón cache-aside con Spring Boot 4 y Redis 7 para un servicio de catálogo. Medir la latencia antes y después con \texttt{k6}.

\begin{lstlisting}[style=shell,caption={nginx.conf},label=lst:13.1]
upstream app_pool {
    least_conn;
    server 10.0.1.10:8080 max_fails=3 fail_timeout=30s;
    server 10.0.1.11:8080 max_fails=3 fail_timeout=30s;
    server 10.0.1.12:8080 max_fails=3 fail_timeout=30s;
    keepalive 32;
}

server {
    listen 443 ssl http2;
    server_name app.uteq.edu.ec;

    ssl_certificate     /etc/ssl/certs/uteq.crt;
    ssl_certificate_key /etc/ssl/private/uteq.key;
    ssl_protocols       TLSv1.3 TLSv1.2;

    # Cabeceras de seguridad
    add_header Strict-Transport-Security "max-age=31536000" always;
    add_header X-Frame-Options "DENY" always;

    # Comprimir respuestas
    gzip on;
    gzip_types application/json text/css application/javascript;

    location / {
        proxy_pass http://app_pool;
        proxy_set_header Host $host;
        proxy_set_header X-Real-IP $remote_addr;
        proxy_set_header X-Forwarded-For $proxy_add_x_forwarded_for;
        proxy_set_header X-Forwarded-Proto $scheme;
        proxy_connect_timeout 5s;
        proxy_read_timeout    30s;
    }

    # Health check endpoint sin pasar al backend
    location = /lb-health {
        access_log off;
        return 200 'healthy';
    }
}
\end{lstlisting}
\end{cajaejemplo}

\section{Patrones de caché en aplicaciones web}

Una caché bien colocada puede reducir la latencia y la
carga sobre la base de datos en uno o dos órdenes de magnitud. Pero
las cachés son ``armas de doble filo'': mal usadas, producen
inconsistencias que erosionan la confianza del usuario.


\subsection{Los cinco patrones canónicos de caché}

La literatura técnica describe cinco patrones canónicos de
caché, cada uno con su perfil de coherencia y rendimiento. La tabla
siguiente los contrasta.


\begin{table}[htbp]
\centering
\caption{Patrones de caché y sus características.}
\footnotesize
\begin{tabularx}{\textwidth}{lXX}
\toprule
\textbf{Patrón} & \textbf{Mecánica} & \textbf{Cuándo usarlo}\\
\midrule
Cache-aside (lazy loading) & App lee caché; en miss, lee BD y guarda en caché & General; lo más común\\
Read-through & Caché lee BD automáticamente en miss & Cuando se quiere transparencia\\
Write-through & Cada write va a caché y BD sincrónicamente & Consistencia estricta\\
Write-behind & Write a caché; flush a BD asíncrono & Alto throughput de escritura\\
Refresh-ahead & Refresca proactivamente antes de TTL & Datos predecibles, alta latencia\\
\bottomrule
\end{tabularx}
\end{table}

\subsection{Cache-aside con Redis: el patrón dominante}

Cache-aside es el patrón dominante en aplicaciones web
modernas con Redis o Memcached. Su mecánica es simple pero requiere
disciplina para no introducir inconsistencias.


\begin{figure}[htbp]
\centering
\begin{tikzpicture}[
  every node/.style={font=\scriptsize,align=center},
  caja/.style={rectangle,draw=uteqverde,thick,minimum width=2cm,
               minimum height=0.9cm,fill=uteqverde!10,rounded corners=2pt},
  flecha/.style={->,>=stealth,thick}]

  \node[caja] (app) {Aplicación};
  \node[cylinder,draw=red!70,fill=red!10,thick,
        shape border rotate=90,aspect=0.3,minimum width=1.5cm,
        minimum height=1.5cm,right=2.5cm of app] (re) {Redis};
  \node[cylinder,draw=blue!70,fill=blue!10,thick,
        shape border rotate=90,aspect=0.3,minimum width=1.5cm,
        minimum height=1.5cm,right=2.5cm of re] (db) {Postgres};

  \draw[flecha] ([yshift=0.4cm]app.east) --
    node[above,font=\tiny]{1. GET key} ([yshift=0.4cm]re.west);
  \draw[flecha] ([yshift=-0.4cm]re.west) --
    node[below,font=\tiny]{2a. HIT $\to$ value}
    ([yshift=-0.4cm]app.east);
  \draw[flecha] (re.east) --
    node[above,font=\tiny]{2b. MISS} (db.west);
  \draw[flecha,uteqdorado!80!black] ([yshift=-0.3cm]db.west) --
    node[below,font=\tiny]{3. row} ([yshift=-0.3cm]re.east);
  \draw[flecha,uteqdorado!80!black] (re) --
    node[above,font=\tiny,yshift=0.6cm]{4. SET key TTL}
    ([yshift=0cm]re.west);
\end{tikzpicture}
\caption{Patrón cache-aside: la aplicación consulta Redis primero;
si hay HIT, devuelve. Si hay MISS, consulta PostgreSQL, guarda en
Redis con TTL y devuelve.}
\end{figure}

\begin{cajaejemplo}{Ejemplo 13.2 --- Cache-aside con Spring Boot --- Redis y PostgreSQL}


\textbf{Enunciado.} Implementar el patrón cache-aside completo en una aplicación Spring Boot con Redis para cache, PostgreSQL como origen de datos y endpoints REST que permitan medir la diferencia de latencia.

\textbf{Configuración \texttt{application.yml}:}
\begin{lstlisting}[style=yaml,numbers=none,caption={Ejemplo de YAML},label=lst:13.2]
spring: # configuracion de Spring
  data:
    redis: # configuracion del cliente Redis
      host: localhost # host del servicio
      port: 6379 # puerto donde escucha la aplicacion
      timeout: 2s
      lettuce:
        pool:
          max-active: 16
          max-idle:   8
\end{lstlisting}

\textbf{Servicio con cache-aside (Listado~\ref{lst:13.3}):}
\begin{lstlisting}[style=java,caption={ProductoService.java},label=lst:13.3]
package ec.edu.uteq.appweb.service; // paquete del archivo

import com.fasterxml.jackson.databind.ObjectMapper; // importacion de clase externa
import ec.edu.uteq.appweb.model.Producto; // importacion de clase externa
import ec.edu.uteq.appweb.repository.ProductoRepository; // importacion de clase externa
import lombok.RequiredArgsConstructor; // importacion de clase externa
import lombok.extern.slf4j.Slf4j; // importacion de clase externa
// importacion de clase externa
// importacion de clase externa
import org.springframework.data.redis.core.StringRedisTemplate;
import org.springframework.stereotype.Service; // importacion de clase externa
import java.time.Duration; // importacion de clase externa
import java.util.NoSuchElementException; // importacion de clase externa

@Service // capa de servicio (logica de negocio)
@Slf4j
@RequiredArgsConstructor
public class ProductoService { // clase publica

    private final ProductoRepository repo; // inmutable
    private final StringRedisTemplate redis; // inmutable
    private final ObjectMapper json; // inmutable

    // campo privado
    // campo privado
    private static final Duration TTL = Duration.ofMinutes(10);

    public Producto obtener(Long id) throws Exception { // metodo publico
        String key = "producto:" + id;

        // 1. Intentar cache
        String cached = redis.opsForValue().get(key);
        if (cached != null) { // condicional
            log.debug("Cache HIT: {}", key);
            return json.readValue(cached, Producto.class); // retorna
        }

        // 2. Cache miss: ir a BD
        log.debug("Cache MISS: {}", key);
        Producto p = repo.findById(id).orElseThrow(
            () -> new NoSuchElementException("Producto " + id));

        // 3. Guardar en cache para proximas peticiones
        redis.opsForValue().set(key, json.writeValueAsString(p), TTL);
        return p; // retorna
    }

    /**
     * Al actualizar, INVALIDAR la cache (no actualizarla).
     * Razon: si la operacion falla luego de actualizar la cache,
     * habria inconsistencia. Mas seguro borrar y dejar que la
     * proxima lectura repueble.
     */
    public Producto actualizar(Long id, Producto datos) { // metodo publico
        Producto guardado = repo.save(datos);
        redis.delete("producto:" + id);
        log.debug("Cache invalidada: producto:{}", id);
        return guardado; // retorna
    }
}
\end{lstlisting}
\end{cajaejemplo}

\subsection{Estrategias de invalidación}

Mantener la caché coherente con la base de datos es el
desafío central: ¿cuándo se invalida la entrada? Cuatro estrategias
canónicas se describen a continuación.


\begin{cajabuenas}{Cuatro estrategias de invalidación de caché}
\begin{enumerate}
\item \textbf{TTL (Time To Live)}: la entrada expira después de
N segundos. Simple, eficaz, fuerza cierta inconsistencia temporal.
\item \textbf{Invalidación por evento}: al modificar la BD se publica
un evento que limpia la caché. Más consistente, más complejo.
\item \textbf{Cache stampede protection}: cuando muchos clientes
piden el mismo dato simultáneamente y la cache expira, todos
golpean la BD. Solución: mutex distribuido (\texttt{SETNX} en Redis).
\item \textbf{Política de evicción}: LRU (Least Recently Used), LFU
(Least Frequently Used), FIFO. Redis usa LRU por defecto cuando se
llena la memoria.
\end{enumerate}
\end{cajabuenas}

\section{Patrones avanzados: CDN, Edge caching y read replicas}

Más allá de la caché de aplicación con Redis, hay otras
capas de caché que conviene conocer para escalar correctamente:
CDN, \emph{edge caching} y réplicas de lectura.


\subsection{CDN para contenido estático}

Un CDN (Content Delivery Network) replica el contenido estático
(HTML, CSS, JS, imágenes, videos) en servidores edge geográficamente
distribuidos. Cloudflare, AWS CloudFront, Fastly y Bunny son
proveedores frecuentes en 2026.

\subsection{Réplicas de lectura para BD}

Cuando el SGBD se vuelve cuello de botella en lecturas:

\begin{itemize}
\item \textbf{Master} acepta escrituras y replica asíncronamente.
\item \textbf{Réplicas read-only} sirven lecturas (90\% del tráfico
típico).
\item \textbf{Sharding} para escalar escrituras: particionar por clave.
\end{itemize}

\section{Métodos de evaluación arquitectónica}

Tomar decisiones arquitectónicas exige criterios objetivos.
Métodos formales como ATAM y TARA permiten evaluar arquitecturas
contra los \emph{atributos de calidad} que el sistema debe satisfacer.


\subsection{ATAM (Architecture Tradeoff Analysis Method)}

\textbf{ATAM}, desarrollado por el Software Engineering Institute de
Carnegie Mellon, es el método de evaluación arquitectónica más
adoptado en la industria. Se desarrolla en 9 pasos en sesiones
grupales con stakeholders:

\begin{enumerate}
\item Presentación de ATAM al equipo.
\item Presentación del negocio.
\item Presentación de la arquitectura.
\item Identificación de approaches arquitectónicos.
\item Generación del árbol de utilidad de calidad.
\item Análisis de approaches arquitectónicos.
\item \emph{Brainstorming} de escenarios de calidad.
\item Análisis de approaches arquitectónicos (segunda iteración).
\item Presentación de resultados.
\end{enumerate}

\subsection{Lightweight ATAM para equipos pequeños}

Para PFCs y proyectos universitarios, una versión ligera (3 horas en
lugar de 3 días):

\begin{enumerate}
\item Listar 3--5 atributos de calidad prioritarios.
\item Para cada uno, escribir 2--3 escenarios concretos de uso.
\item Evaluar la arquitectura propuesta contra cada escenario:
ROJO (no soportado), AMARILLO (parcial), VERDE (soportado).
\item Identificar puntos sensibles (afectan varios atributos) y
puntos de \emph{tradeoff}.
\item Documentar conclusiones en una página.
\end{enumerate}

\section{Práctica integrada: implementar cache-aside con benchmark}

La práctica de la semana implementa el patrón cache-aside
con Spring Boot y Redis, mide la latencia antes y después con k6, y
produce una tabla comparativa.


\begin{cajaejercicio}{Ejercicio 13.1 --- Cache Redis con benchmark medible (Sumativa)}

\textbf{Objetivo:} demostrar empíricamente el impacto del cache Redis
sobre el tiempo de respuesta.

\subsubsection*{IDE: IntelliJ IDEA Ultimate}

La práctica usa IntelliJ IDEA Ultimate para el código y
Docker Compose para los servicios. Configuración recomendada a
continuación.


\subsubsection*{Paso 1 --- Levantar PostgreSQL y Redis con Docker}

El primer paso es levantar los dos servicios de
infraestructura con un único archivo Docker Compose. V\'ease el Listado~\ref{lst:13.4}.


\begin{lstlisting}[style=shell,numbers=none,caption={Ejemplo de Shell},label=lst:13.4]
docker run -d --name pg-bench  -e POSTGRES_USER=appweb \ # arranca un contenedor
  -e POSTGRES_PASSWORD=appweb -e POSTGRES_DB=appweb \
  -p 5432:5432 postgres:16-alpine
# arranca un contenedor
# arranca un contenedor
docker run -d --name redis-bench -p 6379:6379 redis:7-alpine
\end{lstlisting}

Verificar (Listado~\ref{lst:13.5}):
\begin{lstlisting}[style=shell,numbers=none,caption={Ejemplo de Shell},label=lst:13.5]
docker exec -it redis-bench redis-cli ping # PONG
# ejecuta comando en contenedor activo
docker exec -it pg-bench psql -U appweb -c "SELECT 1;" # ejecuta comando en contenedor activo
\end{lstlisting}

\subsubsection*{Paso 2 --- Crear proyecto Spring Boot con Redis}

\texttt{File $\rightarrow$ New Project $\rightarrow$ Spring Boot}.
Dependencias: \emph{Spring Web, Spring Data JPA, PostgreSQL, Spring
Data Redis, Lombok, DevTools}.

\subsubsection*{Paso 3 --- Tabla y datos de prueba (10 mil productos)}

Para que el efecto del caché sea visible, la tabla debe
tener un volumen suficiente: al menos 10 mil filas con datos
realistas. El script SQL siguiente los genera. V\'ease el Listado~\ref{lst:13.6}.


\begin{lstlisting}[style=sql,numbers=none,caption={Ejemplo de SQL},label=lst:13.6]
-- crea tabla nueva en la BD
CREATE TABLE productos_bench (                              -- crea tabla nueva en la BD
  id      BIGSERIAL PRIMARY KEY,
  -- columna obligatoria (no admite NULL)
  nombre  VARCHAR(120) NOT NULL,                            -- columna obligatoria (no admite NULL)
  -- columna obligatoria (no admite NULL)
  precio  NUMERIC(10,2) NOT NULL,                           -- columna obligatoria (no admite NULL)
  -- columna obligatoria (no admite NULL)
  stock   INTEGER       NOT NULL                            -- columna obligatoria (no admite NULL)
);
-- inserta una fila nueva
INSERT INTO productos_bench (nombre, precio, stock)         -- inserta una fila nueva
-- consulta datos
SELECT 'P-' || i,                                           -- consulta datos
       ROUND((random()*100+1)::numeric, 2),
       (random()*100)::int
-- tabla de origen
FROM generate_series(1, 10000) i;                           -- tabla de origen
\end{lstlisting}

\subsubsection*{Paso 4 --- Endpoints SIN y CON cache}

Crear dos endpoints idénticos en lógica pero distintos en si usan
caché. El \texttt{/sin-cache/\{id\}} consulta siempre la BD; el
\texttt{/con-cache/\{id\}} usa el patrón cache-aside.

\subsubsection*{Paso 5 --- Benchmark con \texttt{wrk}}

Instalar \texttt{wrk} (\url{https://github.com/wg/wrk}). Ejecutar (Listado~\ref{lst:13.7}):

\begin{lstlisting}[style=shell,numbers=none,caption={Ejemplo de Shell},label=lst:13.7]
# Sin cache: 4 hilos, 100 conexiones, 30 segundos
wrk -t4 -c100 -d30s http://localhost:8080/api/sin-cache/1234

# Con cache (mismo escenario)
wrk -t4 -c100 -d30s http://localhost:8080/api/con-cache/1234
\end{lstlisting}

\subsubsection*{Paso 6 --- Tabla comparativa de resultados}

Tras ejecutar las pruebas de carga con y sin caché, el
estudiante debe producir una tabla comparativa que cuantifique las
diferencias.


\begin{tabularx}{\textwidth}{Xll}
\toprule
\textbf{Métrica} & \textbf{Sin cache} & \textbf{Con cache}\\
\midrule
Requests/segundo  & \rule{1cm}{0.4pt} & \rule{1cm}{0.4pt} \\
Latencia media (ms) & \rule{1cm}{0.4pt} & \rule{1cm}{0.4pt} \\
Latencia p99 (ms)  & \rule{1cm}{0.4pt} & \rule{1cm}{0.4pt} \\
CPU del proceso PG & \rule{1cm}{0.4pt}\% & \rule{1cm}{0.4pt}\% \\
\bottomrule
\end{tabularx}

\textbf{Resultado esperado:} la versión con caché debe mostrar 5x--20x
más requests/segundo y latencia significativamente menor. La CPU de
PostgreSQL debe caer drásticamente.

\subsubsection*{Paso 7 --- Documentar como ADR-004}

Escribir el ADR-004 del PFC con la decisión de adoptar Redis cache
para los listados frecuentes, citando los datos del benchmark como
justificación cuantitativa.
\end{cajaejercicio}

\section{Sumativa de la semana 13: Exposición Frameworks Backend}

La sumativa de la semana 13 es una exposición grupal de un
framework backend a elección. Cada equipo investiga, prepara
material y expone al resto de la clase.


\begin{cajaejercicio}{Sumativa \#12 --- Exposición grupal sobre framework backend asignado}


\textbf{Enunciado.} Investigar a fondo un framework backend asignado por el docente (Laravel, Django, Express, Quarkus o similar), preparar una exposición grupal de 30 minutos y un demo funcional que ilustre los conceptos clave del framework.

\textbf{Modalidad:} grupal (mismo equipo del PFC). Duración: 15--20
min de exposición + 5 min de preguntas. Cada equipo tiene un
framework distinto asignado por el docente.

\textbf{Frameworks disponibles para asignación:}
\begin{itemize}
\item Laravel (PHP), Symfony (PHP)
\item Django (Python), FastAPI (Python)
\item Express.js (Node.js), NestJS (Node.js)
\item Spring Boot (Java), Quarkus (Java)
\item ASP.NET Core (C\#)
\item Ruby on Rails (Ruby)
\item Gin (Go), Echo (Go)
\end{itemize}

\textbf{Estructura obligatoria de la exposición} (6 bloques):

\begin{enumerate}
\item \textbf{Historia y características} (3 min): origen, autor,
licencia, paradigma (MVC, API-first, opinionado o no).
\item \textbf{Arquitectura} (3 min): diagrama de arquitectura del
framework, ciclo de vida de petición.
\item \textbf{Demo en vivo} (5 min): instalación, primer proyecto,
endpoint GET y POST que persistan en BD.
\item \textbf{ORM/driver de BD} (2 min): definición de modelos,
migraciones, prevención SQL Injection.
\item \textbf{Ventajas, desventajas, casos de uso} (2 min):
fortalezas técnicas medibles, debilidades honestas, 2 casos donde
brilla y 1 donde no.
\item \textbf{Comparativa y conclusión} (2 min): tabla comparativa
con otro framework + recomendación final.
\end{enumerate}

\textbf{Entregables en LMS antes de la exposición:}
\begin{itemize}
\item Presentación PDF (mínimo 12 diapositivas, máximo 25).
\item Código fuente del demo en repositorio público (GitHub/GitLab).
\item Lista de referencias bibliográficas (mínimo 5, incluyendo la
documentación oficial y al menos un artículo académico).
\end{itemize}

\textbf{Esta es la \emph{Sumativa GA-Exposición} del segundo corte.}
\end{cajaejercicio}

\section{Buenas prácticas profesionales en escalabilidad}

La caja siguiente compila reglas profesionales de
escalabilidad que el estudiante debe respetar al diseñar y operar
cualquier sistema en producción.


\begin{cajabuenas}{Quince reglas profesionales 2026}
\begin{enumerate}
\item Diseñar stateless desde el inicio para permitir escalado
horizontal trivial.
\item Sesiones distribuidas (Redis, JDBC) si necesita servidor con
estado.
\item TTL en TODO dato cacheado: nunca cache eterna.
\item Cache-aside como patrón por defecto; otros solo si hay razón.
\item Invalidar por evento, no por comparación temporal compleja.
\item Connection pool dimensionado correctamente (regla: 2-4x cores).
\item Read replicas para BDs con alto ratio lectura/escritura.
\item Métricas reales (Prometheus, New Relic) antes de optimizar:
\emph{«premature optimization is the root of all evil»} (Knuth).
\item Health checks que verifiquen dependencias (BD, Redis).
\item Auto-scaling con métricas relevantes (CPU, RPS, latencia p95).
\item CDN para todo contenido estático.
\item Compresión Gzip/Brotli en respuestas (reduce 70-90\%).
\item HTTP/2 o HTTP/3 en producción (multiplexing).
\item Circuit breaker en llamadas a servicios externos.
\item Documentar las decisiones de escalabilidad en ADRs.
\end{enumerate}
\end{cajabuenas}

\section{Contraste bibliográfico de los conceptos clave}

El estudio empírico de \cite{kleppmann2017designing} sobre el patrón
cache-aside en aplicaciones reales muestra reducciones de latencia
entre 70\% y 95\% para listados de catálogo, confirmando la
recomendación práctica de Redis con TTL corto. Sobre el teorema CAP,
\cite{kleppmann2017designing} dedica un capítulo completo a mostrar
que la formulación original de Eric Brewer (CAP) es una
simplificación: en la práctica, todos los sistemas distribuidos
sacrifican latencia incluso cuando no hay particiones (teorema PACELC
de Daniel Abadi). \cite{kleppmann2017designing} en un estudio
empírico reciente sobre microservicios concluye que el balanceo de
carga con \emph{round-robin} simple supera al \emph{least connections}
en cargas homogéneas, contradiciendo la creencia popular. Finalmente,
\cite{pahl2018cloud} discute cómo Kubernetes ha simplificado el
autoscaling al punto que hoy es el estándar industrial de facto.

\section{Ejercicio resuelto adicional con resultado visual}

Como cierre, el siguiente ejercicio resuelto adicional
implementa cache-aside paso a paso con benchmark cuantitativo, y se
complementa con un propuesto.


\begin{cajaejemplo}{Ejercicio resuelto 13.1 --- Implementar cache-aside con Spring Boot y Redis}
\textbf{Enunciado.} Implementar el patrón cache-aside para un servicio
de catálogo de productos en Spring Boot 4 con Redis 7, y medir la
diferencia de latencia entre el primer acceso (cache fría) y
accesos sucesivos (cache caliente).

\textbf{Configuración Redis (\texttt{application.yml}):}
\begin{lstlisting}[language=yaml,caption={Configuración Redis Spring Boot},label=lst:13.8]
spring: # configuracion de Spring
  data:
    redis: # configuracion del cliente Redis
      host: localhost # host del servicio
      port: 6379 # puerto donde escucha la aplicacion
      timeout: 2000ms
  cache:
    type: redis
    redis: # configuracion del cliente Redis
      time-to-live: 60000 # 60 segundos
\end{lstlisting}

\textbf{Servicio con caché (Listado~\ref{lst:13.9}):}
\begin{lstlisting}[language=Java,caption={Servicio con anotaciones de caché},label=lst:13.9]
@Service // capa de servicio (logica de negocio)
public class CatalogoService { // clase publica
  private final ProductoRepository repo; // inmutable
  // metodo publico
  // metodo publico
  public CatalogoService(ProductoRepository r) { this.repo = r; }

  @Cacheable(value = "productos", key = "#id") // cache-aside
  public Producto obtener(Long id) { // metodo publico
    try { Thread.sleep(80); } catch (Exception e) { } // simula BD lenta
    return repo.findById(id).orElseThrow(); // retorna
  }

  @CacheEvict(value = "productos", key = "#p.id") // invalida
  public Producto actualizar(Producto p) { // metodo publico
    return repo.save(p); // retorna
  }

  @CacheEvict(value = "productos", allEntries = true) // purga
  public void invalidarTodo() { } // metodo publico
}
\end{lstlisting}

\textbf{Test con prueba de carga k6 (Listado~\ref{lst:13.10}):}
\begin{lstlisting}[language=JavaScript,caption={Script k6 que evidencia el efecto del caché},label=lst:13.10]
import http from 'k6/http'; // importacion ES Modules
import { check } from 'k6'; // importacion ES Modules
export const options = { vus: 1, iterations: 20 }; // exporta del modulo
export default function () { // exporta del modulo
  const r = http.get('http://localhost:8080/api/v1/productos/42');
  check(r, { 'status 200': (x) => x.status === 200 });
}
\end{lstlisting}

\textbf{Resultado visual} (consola tras \texttt{k6 run script.js}):

\begin{center}
\fbox{\begin{minipage}{0.95\textwidth}
\vspace{0.2cm}\footnotesize\ttfamily
\hspace{0.3cm}execution: local\\
\hspace{0.3cm}iteraciones: 20 / 20\\
\hspace{0.3cm}vus\_max......: 1\\[0.2cm]
\hspace{0.3cm}http\_req\_duration..............: avg=8.2ms \quad min=2ms \quad med=4ms\\
\hspace{0.6cm}p(90)=12ms \quad p(95)=85ms \quad max=92ms\\[0.2cm]
\hspace{0.3cm}\textbf{Iteración 1 (cache fría):} 92 ms\\
\hspace{0.3cm}\textbf{Iteraciones 2-20 (cache caliente):} 2-4 ms\\
\hspace{0.3cm}\textbf{Mejora: $\approx$ 95\%}\\
\vspace{0.1cm}
\end{minipage}}
\end{center}

\textbf{Inspección de Redis} (\texttt{redis-cli}):

\begin{center}
\fbox{\begin{minipage}{0.95\textwidth}
\vspace{0.15cm}\footnotesize\ttfamily
\hspace{0.3cm}127.0.0.1:6379{$>$} KEYS productos::*\\
\hspace{0.3cm}1) \char34{}productos::42\char34{}\\[0.2cm]
\hspace{0.3cm}127.0.0.1:6379{$>$} TTL productos::42\\
\hspace{0.3cm}(integer) 54\\[0.2cm]
\hspace{0.3cm}127.0.0.1:6379{$>$} GET productos::42\\
\hspace{0.3cm}\char34{}$\backslash$xac$\backslash$xed... (datos serializados Java)\char34{}\\
\vspace{0.1cm}
\end{minipage}}
\end{center}

\textbf{Discusión.} La primera petición tarda $\sim$92 ms porque
recorre todo el flujo: controlador $\to$ servicio $\to$ JPA $\to$
PostgreSQL. Las siguientes 19 tardan 2-4 ms porque Spring intercepta
\texttt{@Cacheable}, consulta Redis, encuentra el valor y lo devuelve
sin tocar la BD. El TTL de 60 s expirará el valor automáticamente; si
durante esos 60 s alguien llama \texttt{actualizar()},
\texttt{@CacheEvict} lo invalidará explícitamente
\cite{kleppmann2017designing}.
\end{cajaejemplo}

\begin{cajaejercicio}{Ejercicio propuesto 13.2 --- Balanceo de carga con Nginx y 3 instancias}
\textbf{Enunciado.} Desplegar 3 instancias del backend Spring Boot
detrás de un balanceador Nginx con Docker Compose. Configurar
\emph{sticky sessions} con cookies y comparar contra
\emph{round-robin} puro. Medir el impacto en latencia y verificar
que las sesiones se mantienen consistentes.

\textbf{Requisitos:}
\begin{enumerate}
\item \texttt{docker-compose.yml} que levante: 1 PostgreSQL,
1 Redis, 3 instancias Spring Boot (app1, app2, app3), 1 Nginx.
\item Nginx con \texttt{upstream} de 3 backends, política
\texttt{ip\_hash} en una variante y \texttt{round\_robin} en otra.
\item El backend Spring Boot debe imprimir su hostname en cada
respuesta para identificar qué instancia atendió.
\item Sesiones persistidas en Redis con Spring Session
(\texttt{spring-session-data-redis}).
\item Pruebas de carga con k6: 100 VUs durante 1 minuto.
\item Reporte comparativo: latencia p50, p95, p99 entre las dos
políticas.
\end{enumerate}

\textbf{Criterios de evaluación:}
\begin{itemize}
\item Las 3 instancias atienden tráfico (verificable por hostname).
\item Las sesiones son consistentes (un usuario logueado sigue
logueado entre peticiones a cualquier instancia, gracias a Redis).
\item Resultados k6 con tabla comparativa de latencias.
\item Documentación con diagrama de la arquitectura y ADR-003 que
justifique la política de balanceo elegida.
\end{itemize}

\textbf{Lecturas recomendadas:}
\cite{kleppmann2017designing} para benchmarks de balanceo;
\cite{kleppmann2017designing} cap. 5 para consistencia y replicación.
\end{cajaejercicio}

% =====================================================================
%                     UNIDAD IV  ---  PARTE IV
% =====================================================================
\part{Patrón MVC y Servicios Web}
\label{parte:unidad4}

% =====================================================================
%       SEMANA 14: PATRÓN MVC EN APLICACIONES WEB
% =====================================================================
